\documentclass[11pt]{article}
\usepackage[letterpaper,margin=1in]{geometry}
\usepackage[T1]{fontenc}
\usepackage{lmodern}
\usepackage{microtype}
\usepackage{amsmath,amssymb,amsthm,mathtools}
\usepackage{booktabs,longtable,tabularx,array}
\newcolumntype{Y}[1]{>{\raggedright\arraybackslash}p{#1}}
\usepackage{enumitem}
\usepackage{xcolor}
\usepackage{listings}
\usepackage{fancyhdr}
\usepackage{hyperref}
\usepackage{xurl}
\usepackage{cleveref}
\usepackage{graphicx}
\usepackage{tikz}
\usetikzlibrary{arrows.meta,positioning,calc,fit}

\definecolor{codebg}{RGB}{247,247,247}
\definecolor{codeframe}{RGB}{210,210,210}
\definecolor{linkblue}{RGB}{24,74,140}
\hypersetup{colorlinks=true,linkcolor=linkblue,citecolor=linkblue,urlcolor=linkblue,pdftitle={FIBER-iPRA},pdfauthor={OpenAI}}
\lstset{basicstyle=\ttfamily\small,backgroundcolor=\color{codebg},frame=single,rulecolor=\color{codeframe},breaklines=true,columns=fullflexible,keepspaces=true,showstringspaces=false}
\pagestyle{fancy}
\fancyhf{}
\lhead{FIBER-iPRA}
\rhead{16 November 2026}
\cfoot{\thepage}
\setlength{\headheight}{14pt}
\setlist{itemsep=2pt,topsep=4pt}

\newtheorem{theorem}{Theorem}[section]
\newtheorem{lemma}[theorem]{Lemma}
\newtheorem{proposition}[theorem]{Proposition}
\newtheorem{corollary}[theorem]{Corollary}
\newtheorem{definition}[theorem]{Definition}
\newtheorem{principle}[theorem]{Principle}
\newtheorem{warning}[theorem]{Warning}
\newtheorem{example}[theorem]{Example}
\newtheorem{observation}[theorem]{Observation}
\theoremstyle{remark}
\newtheorem{remark}[theorem]{Remark}

\newcommand{\ipra}{\mathrm{iPRA}}
\newcommand{\PRAAx}{\mathsf{PRAAx}}
\newcommand{\Prf}{\mathsf{Prf}}
\newcommand{\Nat}{\mathbb N}
\newcommand{\Bool}{\{0,1\}}
\newcommand{\qf}{\mathsf{QF}}
\newcommand{\cert}{\mathsf{Cert}}
\newcommand{\checkb}{\mathsf{check}}
\newcommand{\obs}{\mathsf{obs}}
\newcommand{\res}{\mathsf{Res}}
\newcommand{\dem}{\mathsf{dem}}
\newcommand{\code}{\ulcorner\!\cdot\!\urcorner}
\newcommand{\seq}{\mathsf{Seq}}
\newcommand{\true}{\mathsf{tt}}
\newcommand{\false}{\mathsf{ff}}
\newcommand{\leqprf}[2]{#1\;\vdash_{\PRAAx}\;#2}

\title{\textbf{FIBER-iPRA}\\[3pt]
\large A Fibered Interface Calculus for Compiling Finitary Mathematics into Intuitionistic Primitive Recursive Arithmetic}
\author{Research report prepared for the ProveIt repository}
\date{16 November 2026\\Repository baseline: \texttt{58627ed89814ba122e1f0f012f1aba103cda267b}}

\begin{document}
\maketitle

\begin{abstract}
The ProveIt repository contains a checked ordinary presentation of intuitionistic primitive recursive arithmetic, a substantial checked arithmetic library, a bridge to the Foundation presentation, and more than one hundred research packets that reformulate mathematics in finite quantitative terms.  The recurring gap is no longer the absence of isolated encodings.  It is the absence of one compositional interface that distinguishes object-theory derivations from host computations, controls the scope of finite certificates, turns bounded quantifiers into quantifier-free loop invariants, and records genuinely missing effective information as explicit residual assumptions.

This report develops \emph{FIBER-iPRA} (Finite-query Interfaces with Bounded Evidence and Realizers).  A mathematical object is represented by a family of finite certificate spaces indexed by a query.  A morphism carries an explicit finite demand map, a certificate transformer, an observation map, and a derivation in the literal judgment \(\Prf\;\PRAAx\;\Gamma\;\varphi\).  Local certificates are eliminated by existential elimination; least bounded search is used only when a public result must name a canonical code.  Composition is proved by substitution and simultaneous cut.  Cyclic resource requirements are resolved by a finite monotone closure theorem.  Fixed bounded quantifier alternations are handled by finite strategy certificates.  Missing constants, cutoffs, moduli, and representatives are exposed as residual producer formulas, with a precise hierarchy of cases in which the residual can be eliminated inside \(\ipra\).

The framework is applied to several difficulties reported in the repository.  We give constructive proofs of a dyadic adjacent-motion theorem that upgrades compatible approximants to a Cauchy name, a finite-net lifting theorem replacing abstract compactness, guarded exact branch selection, stable cluster-valued observations at collisions, finite master-certificate formation without hidden choice, and bounded-game realization.  We also prove two useful impossibility results: there is no primitive-recursive universal evaluator for all primitive-recursive functions, and no finite-use exact sign or branch selector can work on arbitrary Cauchy names without a separation guard.  A dependency-ordered Lean plan and source-oriented Lean modules use the current \texttt{ConciseIPRA}, \texttt{IPRADerivations}, and \texttt{IPRABridge} interfaces.  No new kernel-checked theorem is claimed by this archive; the report states exactly which declarations must be completed and how each reduces to existing checked infrastructure.
\end{abstract}

\paragraph{Verification status.}
The mathematical arguments in this report are human proofs.  The Python reference implementation and tests exercise only finite algorithms.  The Lean files are an integration blueprint tied to the current repository API; Lean was unavailable in the artifact runtime, so those files were not elaborated.  A theorem is counted as formally established only when the repository contains a compiled declaration of literal type \texttt{Prf PRAAx ...} and its axiom audit passes.

\tableofcontents
\newpage

\section{Repository audit and the actual problem}
\label{sec:audit}

\subsection{Audited baseline}

The report was prepared against the public \texttt{main} branch of
\texttt{VladimirReshetnikov/ProveIt} at commit
\texttt{58627ed89814ba122e1f0f012f1aba10\allowbreak 3cda267b}.  At that
snapshot the \path{Logic/PrimitiveRecursiveArithmetic} directory has four logically
distinct layers.

\begin{longtable}{p{0.20\linewidth}p{0.47\linewidth}p{0.24\linewidth}}
\toprule
Layer & Role & Status at the audited snapshot\\
\midrule
\endfirsthead
\toprule
Layer & Role & Status at the audited snapshot\\
\midrule
\endhead
\texttt{ConciseIPRA} & Ordinary first-order \(\ipra\), its small native scalar core, and a bi-interpretation.  The portable judgment is \(\Gamma\vdash_{\PRAAx}\varphi\), written in Lean as \texttt{Prf PRAAx Gamma phi}. & Lean library checked by the kernel.\\
\texttt{Derivations} & Reusable arithmetic and theorem-specific derivations: semiring laws, order, monus, pairing, bounded search, division, gcd, B\'ezout, finite sums, descent, linear algebra, and many completed packet theorems. & Lean library checked by the kernel.\\
\texttt{Bridge} & Translation between ordinary ConciseIPRA and the Foundation \(\ipra\) used by several computability projects, including open-sequent transport and arithmetic symbol comparison. & Lean library checked by the kernel.\\
\texttt{docs/proofs} & 107 delivered packets in algebra, analysis, number theory, geometry, probability, topology, and foundations.  They contain human arguments, executable checks, and many Lean blueprints. & Filed as delivered; the directory README explicitly says that packets were not rerun or compiled on intake.\\
\bottomrule
\end{longtable}

The distinction is decisive.  A theorem of Lean about \texttt{Nat}, a Python
calculation, a verified polynomial identity in the host, or a structure whose
field asks for a proof is not yet an \(\ipra\) theorem.  The acceptance test is
literal: after all adapters have been expanded, the endpoint must have type
\[
        \Gamma\;\vdash_{\PRAAx}\;\varphi.
\]
The repository's checked \texttt{Derivations} library meets this test.  Most
new proof packets do not, and their own status notes are generally careful
about that fact.

\subsection{What the existing foundation packets already solve}

The most recent foundation packets converge on a strong common picture.
Finite-query normal form packs all evidence for a fixed query into one bounded
certificate.  KIFC-R and FINER add explicit residual producers and finite
demand closure.  CBIP isolates bounded interaction programs and finite
strategy certificates.  Scoped Certificate Realization observes that most
certificates should remain local and be eliminated by existential elimination.
The observation-algebra packets replace discontinuous labels by stable
invariants or guarded labels.  These are genuine advances, and the present
work deliberately reuses them rather than renaming them.

Their recurring unresolved obligations are also remarkably stable:

\begin{enumerate}
\item a host-level finite algorithm has not been turned into an object-level
      derivation;
\item a bounded formula is placed inside an induction predicate, making the
      predicate non-quantifier-free;
\item finite witnesses are chosen independently without a compatibility proof;
\item a missing modulus, separation bound, height cutoff, effective constant,
      or representative is hidden in host data;
\item a noncanonical witness escapes into a public conclusion;
\item an infinitary object (a completion, compact space, quotient, sheaf, or
      probability space) is invoked where the actual proof uses only finitely
      many observations;
\item metastability or pointwise approximation is silently promoted to a
      coherent name or a global rate;
\item a uniform evaluator or exact selector is demanded even though the
      proposed interface cannot exist.
\end{enumerate}

The current repository already contains the arithmetic ingredients needed for
most of the finite part: arbitrary-term defining equations, quantifier-free
induction, substitution, simultaneous cut, order, bounded search, pairing,
simultaneous recursion, bounded sums and products, and finite linear algebra.
The missing shared layer is therefore not a larger arithmetic library.  It is
a disciplined compiler interface connecting a finite mathematical argument to
those checked derivations.

\subsection{Why another framework is justified}

A new framework is useful only if it removes a repeated source of work rather
than adding another vocabulary.  FIBER-iPRA makes three changes.

\paragraph{First, interfaces are query-indexed and typed.}
A ``real number'', ``root'', ``finite cover'', or ``cohomology class'' is not
represented by one global code.  It is represented by what can be verified at
a requested precision, degree, horizon, mesh, or finite diagram.  The
certificate space is allowed to vary with that query.  This makes resource
flow explicit and prevents a proof at one precision from being confused with
a coherent name at every precision.

\paragraph{Second, composition is an object-theory theorem.}
Every reusable node comes with a local derivation.  The compiler specializes
these derivations, substitutes actual terms, and composes them using the
checked \texttt{Prf.subst}, \texttt{Prf.weaken}, \texttt{Prf.exE}, and
\texttt{Prf.trans} operations.  ``The checker is correct'' is not left as an
informal bridge from executable truth to \(\ipra\).

\paragraph{Third, residuals have a calculus.}
A missing producer is neither concealed nor automatically fatal.  It is a
formula in the final sequent.  The report classifies when it can be eliminated
by bounded search, master-certificate formation, finite-net construction,
adjacent-motion estimates, or finite strategy realization.  Residuals that
encode genuinely unavailable effective mathematics remain explicit.

\subsection{Conservative design rule}

FIBER-iPRA does not add axioms to \(\ipra\) and does not propose an internal
universal interpreter.  A fixed external source generator produces fixed
primitive-recursive symbols, formulas, and proof terms.  Lean checks the
resulting derivations.  This is exactly the same trusted-boundary pattern used
successfully by the checked arithmetic library: metatheory may help construct
a proof, but the object theorem is accepted only through the inductive
\texttt{Prf} datatype.

\section{The formal boundary: ordinary iPRA and finite specialization}
\label{sec:boundary}

\subsection{The object theory used in the report}

We use the ordinary presentation already selected by the repository
unification plan.  Its terms are built from a symbol for every primitive
recursive function.  Its formulas use intuitionistic first-order logic with
equality.  The theory \(\PRAAx\) contains the defining equations for primitive
recursive symbols, the successor axioms, and induction for quantifier-free
formulas.  In Lean the central judgment is
\[
       \Prf\;\PRAAx\;\Gamma\;\varphi,
\]
with notation \(\Gamma\vdash_{\PRAAx}\varphi\).

The checked deduction layer provides, among other rules:

\begin{itemize}
\item substitution of arbitrary terms for free variables;
\item weakening of the hypothesis list;
\item simultaneous cut, called \texttt{Prf.trans};
\item universal and existential introduction and elimination;
\item simultaneous equality substitution;
\item quantifier-free induction through \texttt{PRA.qfInd};
\item the convenient open induction wrapper \texttt{LeanProofs.IPRA.ind};
\item instantiation of closed derivations at arbitrary terms through
      \texttt{LeanProofs.IPRA.inst}.
\end{itemize}

The framework never reasons by saying that a primitive recursive function has
the intended standard value and then applying a semantic completeness
principle.  Instead, every use of a program is mediated by its defining
equations and an object derivation.

\subsection{Four levels that must not be conflated}

For a fixed algorithm \(f\), formula \(A\), and certificate \(c\), there are
four different statements.

\begin{enumerate}
\item \textbf{Host computation:} a Python or Lean evaluator returns a value.
\item \textbf{Host theorem:} Lean proves a proposition about the standard
      natural-number interpretation of \(f\).
\item \textbf{Object equation:} \(\ipra\) derives the defining or correctness
      equation for the term denoting \(f\).
\item \textbf{Object conclusion:} \(\ipra\) derives \(A\), perhaps after
      eliminating the certificate and discharging residual assumptions.
\end{enumerate}

Levels 1 and 2 are useful for discovering invariants and regression testing,
but only levels 3 and 4 prove the requested theorem.  FIBER metadata records
all four separately.

\subsection{Why the compiler is specialized}

One might try to internalize the entire framework by adding a
primitive-recursive universal evaluator \(U(e,x)\) for codes of arbitrary
primitive-recursive functions.  This cannot work.

\begin{theorem}[No primitive-recursive universal evaluator]
There is no binary primitive-recursive function \(U\) such that for every
unary primitive-recursive function \(f\) there is an index \(e\) satisfying
\(U(e,x)=f(x)\) for every \(x\).
\end{theorem}

\begin{proof}
If such \(U\) existed, the diagonal function
\[
             g(x)=U(x,x)+1
\]
would be primitive recursive.  Hence \(g(x)=U(e,x)\) for some \(e\).  Setting
\(x=e\) gives \(U(e,e)+1=U(e,e)\), impossible.  The argument is constructive
and uses only closure of primitive recursive functions under composition and
successor.
\end{proof}

Consequently a framework that quantifies internally over all primitive
recursive source programs is asking for the wrong interface.  FIBER uses
\emph{metalevel specialization}: a finite source description is compiled to
specific \(\mathsf{PR}\) terms, specific formulas, and a specific proof term.
For a bounded dynamic trace, the trace checker itself is a fixed primitive
recursive symbol.  The generated proof is then checked by Lean's kernel.

\subsection{Finite-use discipline}

A proof of a fixed finite-query conclusion can inspect only finitely much
source data.  FIBER makes that finite use explicit through a demand map.  This
is not merely an implementation optimization.  It is the logical reason that
objects such as Cauchy reals, analytic functions, compact spaces, and quotient
objects can be handled in weak arithmetic: the theorem does not need the
completed infinitary object; it needs a finite observation and a proof that
this observation is stable under the chosen representation.

The discipline has two forms.

\begin{description}
\item[Static finite use.] The compiler knows a finite list of source queries
      from the target query.
\item[Adaptive finite use.] A bounded interaction program chooses its next
      query from previous answers, but a primitive-recursive fuel or a finite
      monotone closure bounds the complete trace.
\end{description}

Both forms produce one finite master certificate.  No appeal to countable
choice or bar induction is needed for the compiled theorem.

\subsection{Portable versus backend-specific statements}

The canonical portable conclusion is an ordinary \(\Prf\;\PRAAx\) sequent.
When an application is stated in the Foundation syntax, the checked bridge
transports the resulting open or closed sequent.  This separation avoids two
common mistakes:

\begin{enumerate}
\item proving the right mathematics in a nonportable host representation and
      never obtaining an object theorem;
\item rebuilding arithmetic inside every application-specific Foundation
      development.
\end{enumerate}

The framework therefore compiles first to the ordinary calculus, imports the
checked arithmetic library, and translates only at the application boundary.

\section{Fibered finite-query interfaces}
\label{sec:interfaces}

\subsection{Encoded finite data}

All finite collections are represented by natural numbers using a fixed
primitive-recursive sequence coding.  We write
\(\seq(C,i)\) for the \(i\)-th entry and \(\ell(C)\) for the declared length.
A code is accepted only under an explicit bound.  The concrete implementation
may use the pairing and simultaneous-recursion infrastructure already checked
in \texttt{IPRADerivations.Core.Pairing}, or mixed-radix coding when component
bounds are known.  The mathematical statements below are independent of that
choice.

Every Boolean predicate is represented by a primitive-recursive bit.  Thus
\(b(\vec x)=1\) is an atomic equality in the object language.  A checker is
never trusted merely because it executes; it is accompanied by a reflection
or soundness derivation.

\subsection{Interfaces}

\begin{definition}[FIBER interface]
A finite-query interface \(X\), with parameter code \(p\), consists of fixed
primitive-recursive data:
\[
\begin{array}{lll}
Q_X(p,q) &\in \Bool, &\text{well-formed query bit},\\
B_X(p,q) &\in \Nat, &\text{exclusive certificate bound},\\
V_X(p,q,c) &\in \Bool, &\text{certificate verifier},\\
O_X(p,q,c) &\in \Nat, &\text{observation code},\\
E_X(p,q,o,o') &\in \Bool, &\text{observational equivalence},\\
D_X(p,q,o,o') &\in \Nat, &\text{optional quantitative discrepancy}.
\end{array}
\]
It also contains object-theory derivations for verifier soundness,
observational equivalence, and congruence of every exported operation.
\end{definition}

The certificate fiber over \(q\) is the finite set
\[
  \cert_X(p,q)=\{c<B_X(p,q):V_X(p,q,c)=1\}.
\]
The point of the word ``fiber'' is that both the admissible evidence and its
bound may depend on the query.  A degree-\(d\) polynomial calculation, a
precision-\(n\) approximation, and a horizon-\(T\) trace need not share a
single global witness.

\begin{definition}[Validity and totality]
The interface formula
\[
  \mathsf{Valid}_X(p,q,c)\;:\equiv\;c<B_X(p,q)\wedge V_X(p,q,c)=1
\]
is quantifier-free.  The totality formula is
\[
  \mathsf{Total}_X(p,q)\;:\equiv\;
  \exists c\,\mathsf{Valid}_X(p,q,c).
\]
Totality may be a theorem, a hypothesis inherited from another interface, or
an explicit residual producer.
\end{definition}

\begin{definition}[Certificate-independent observation]
An observation is certificate-independent when \(\ipra\) derives
\[
  \mathsf{Valid}_X(p,q,c)\wedge\mathsf{Valid}_X(p,q,c')
  \longrightarrow E_X(p,q,O_X(p,q,c),O_X(p,q,c'))=1.
\]
For a metric interface one may instead derive a discrepancy bound depending
only on the query.
\end{definition}

This weaker notion is often exactly what mathematics supplies.  Different
root isolations may yield the same cluster, different finite covers may yield
the same stable count, and different cocycle representatives may yield the
same finite invariant.  Code equality is neither required nor desirable.

\subsection{Realizers and demand}

\begin{definition}[FIBER realizer]
A realizer \(R:X\to Y\) consists of:

\begin{enumerate}
\item a primitive-recursive demand list \(\dem_R(p,q)\) of \(X\)-queries
      needed to answer the \(Y\)-query \(q\);
\item a primitive-recursive bound for a master input certificate;
\item a master verifier asserting that every demanded query has a valid
      certificate;
\item a certificate transformer \(T_R(p,q,C)\) producing a \(Y\)-certificate;
\item an observation transformer, when the public output is an observation;
\item a literal object derivation of the local soundness sequent;
\item extensionality or stability derivations showing that equivalent input
      observations give equivalent output observations;
\item a list of residual producer formulas not supplied by \(X\).
\end{enumerate}
\end{definition}

The local soundness sequent has the schematic form
\[
  \res_R(p,q),\;
  \bigl\{\mathsf{Valid}_X(p,q_i,c_i):q_i\in\dem_R(p,q)\bigr\}
  \;\vdash_{\PRAAx}\;
  \mathsf{Valid}_Y(p,q,T_R(p,q,C)).
\]
Here \(\res_R(p,q)\) abbreviates a finite list of exact object-language
formulas.  There is no semantic leap from the fact that a host transformer
returns a certificate.

\subsection{Local, public, and residual evidence}

The interface records the \emph{escape class} of each witness.

\begin{description}
\item[Local.] The witness is introduced to prove an intermediate statement and
      eliminated before the node returns.  It is never named by a public term.
\item[Observational.] The code may vary, but all valid codes yield equivalent
      or quantitatively close observations.  The code is eliminated; only the
      observation class escapes.
\item[Canonical.] A downstream theorem genuinely needs one code.  A least
      valid code is selected by bounded search, after totality has been proved.
\item[Residual.] Existence of the witness is not established by the current
      finite library.  Its exact producer formula remains in the final sequent.
\end{description}

This classification prevents a subtle but frequent strengthening.  From
\(\forall q\,\exists c\,V(q,c)\) one may use \(c\) locally for a fixed \(q\).
One may not silently obtain a global function \(q\mapsto c_q\).  FIBER obtains
such a function only when the verifier is bounded and decidable, so checked
least search defines it primitively recursively.

\subsection{Herbrand capsules}

\begin{definition}[Herbrand capsule]
A capsule for a relation \(A(q,w)\) is a bounded list
\(W=\langle w_0,\ldots,w_{k-1}\rangle\), together with a verifier proving
\[
     \bigvee_{i<k} A(q,w_i).
\]
The list, its length, and the finite disjunction are encoded by primitive
recursive data and a fixed proof schedule.
\end{definition}

A capsule is the preferred output when a proof establishes a bounded finite
set of possibilities but no stable reason to select one.  If a downstream
operation is compatible across all valid members, the capsule can be mapped
pointwise and eliminated.  If a canonical member is required and \(A\) is a
reflected bounded predicate, least search selects the first valid member.
Thus capsules isolate exactly where finite choice is harmless and where a
selection theorem is required.

\subsection{Metric and quotient fibers}

Two especially useful interface flavors are built into the framework.

\paragraph{Metric fibers.}
For precision query \(n\), valid certificates yield observations within an
explicit error \(\epsilon(n)\).  Composition adds or otherwise transforms the
errors.  No equality of approximants is expected.

\paragraph{Quotient fibers.}
An equivalence bit \(E_X\) is reflexive, symmetric, transitive, and respected
by exported operations, all by object derivations.  Public statements are
phrased in terms of invariant observations.  A global quotient type need not
be constructed inside \(\ipra\).

These two forms cover many apparent obstacles involving real numbers,
projective directions, unordered roots, eigenspaces, local coordinates,
finite covers, and cohomology classes.

\subsection{A schematic example: rational root clusters}

Let \(p\) code a monic rational polynomial and let the query \(q=(N,r)\)
request information about roots in a rational box at mesh \(2^{-N}\), with
multiplicity cutoff \(r\).  A certificate can contain a list of rational boxes,
Sturm or resultant signs, and multiplicities.  The verifier checks coverage,
disjointness at the requested mesh, and the coefficient identities.  The
observation is the multiset of boxes and multiplicities, or the first \(r\)
elementary symmetric functions.  At a collision there is no exact ordered
root label, but the multiset observation remains stable.  Under an additional
separation residual, the same interface refines to ordered labels.  Nothing in
the representation requires a discontinuous global branch selector.

\section{The finite proof kernel}
\label{sec:kernel}

This section states the reusable lemmas from which the compiler is built.
They are intentionally small.  Each either follows directly from the checked
natural-deduction calculus or reduces to quantifier-free induction and checked
arithmetic.

\subsection{Certificate sealing}

\begin{theorem}[Local certificate sealing]
\label{thm:sealing}
Let \(A(c,\vec x)\) and \(B(\vec x)\) be formulas, with \(c\) absent from
\(B\).  If
\[
 \Gamma\vdash_{\PRAAx}\exists c\,A(c,\vec x)
 \qquad\text{and}\qquad
 A(c,\vec x),\Gamma^+\vdash_{\PRAAx}B(\vec x)^+,
\]
then \(\Gamma\vdash_{\PRAAx}B(\vec x)\).
\end{theorem}

\begin{proof}
This is exactly the existential-elimination constructor \texttt{Prf.exE} of
\texttt{ConciseIPRA.Deduction}.  The shifts in the displayed second premise
express that the fresh certificate variable is in scope only in the branch.
No choice function is created.
\end{proof}

\begin{corollary}[Observational escape]
Suppose every valid certificate \(c\) gives a conclusion \(B(o)\) for
\(o=O(q,c)\), and valid certificates give equivalent observations.  If the
public conclusion is invariant under that equivalence, the certificate can be
sealed and only the invariant conclusion escapes.
\end{corollary}

\begin{proof}
Open the totality witness with \cref{thm:sealing}, derive the invariant
conclusion from its observation, and close the witness.  Equivalence is needed
only when two independently generated certificates are compared.
\end{proof}

\subsection{Simultaneous residual discharge}

\begin{theorem}[Finite residual discharge]
\label{thm:cut}
If \(\Gamma\vdash_{\PRAAx}\varphi\) and for every \(\psi\in\Gamma\) one has
\(\Delta\vdash_{\PRAAx}\psi\), then
\(\Delta\vdash_{\PRAAx}\varphi\).
\end{theorem}

\begin{proof}
This is the checked theorem \texttt{Prf.trans}.  Its proof repeatedly applies
implication introduction and modus ponens, and therefore requires no cut rule
as a new primitive.
\end{proof}

This theorem is the formal meaning of ``plugging producers into a compiled
proof''.  A framework report that lists residual assumptions but does not
produce this sequent has not yet specified composition at the object level.

\subsection{Booleanization of fixed bounded formulas}

For each fixed bounded formula \(A(\vec x)\), FIBER generates a primitive
recursive bit \([A](\vec x)\).  Atomic equalities and inequalities use the
checked equality/order bits.  Connectives use Boolean operations.  A bounded
existential is a finite disjunction fold, and a bounded universal is a finite
conjunction fold.

\begin{theorem}[Fixed-formula Boolean reflection schema]
\label{thm:boolean-reflection}
For every fixed bounded formula \(A(\vec x)\), there is a primitive-recursive
bit \([A]\) and a generated object derivation
\[
  []\vdash_{\PRAAx}\forall\vec x\,
       \bigl([A](\vec x)=1\leftrightarrow A(\vec x)\bigr).
\]
The construction is effective in the external syntax of \(A\); it does not
require an internal universal formula evaluator.
\end{theorem}

\begin{proof}[Human proof and induction accounting]
Proceed by structural recursion on the fixed syntax.

For equality and order atoms, use the existing checked reflection lemmas for
equality tests, monus order, and comparison bits.  Boolean conjunction,
disjunction, and implication are proved by the finite truth tables, each a
fixed collection of arithmetic cases.

For a bounded existential \(\exists i<N\,A(i,\vec x)\), define
\[
 e(0,\vec x)=0,\qquad
 e(k+1,\vec x)=\max(e(k,\vec x),[A](k,\vec x)).
\]
The formula used in induction is not the quantified formula.  It is a
quantifier-free state equation about \(e\), a candidate index, and the
reflected bit.  The forward direction from \(e(N)=1\) obtains the first stage
at which the bit becomes one by checked bounded search, then applies the
inductive reflection theorem for \(A\).  The reverse direction takes an
explicit existential witness \(i<N\), proves by a quantifier-free loop lemma
that its bit contributes to \(e(N)\), and concludes \(e(N)=1\).

For a bounded universal, define the conjunction fold
\[
 a(0,\vec x)=1,\qquad
 a(k+1,\vec x)=\min(a(k,\vec x),[A](k,\vec x)).
\]
From \(a(N)=1\) and an arbitrary \(i<N\), a quantifier-free indexing lemma
proves \([A](i,\vec x)=1\), then reflection yields \(A(i,\vec x)\).  Conversely,
under the universal hypothesis, quantifier-free induction proves the state
formula \(a(k)=1\) for \(k\le N\); the successor step instantiates the
universal hypothesis at \(k\).  The induction predicate itself remains the
atomic equation \(a(k)=1\), optionally guarded by the quantifier-free
inequality \(k\le N\).

All recursion symbols and formulas are fixed before the derivation is
constructed.  Thus no truth predicate and no universal evaluator is used.
\end{proof}

\begin{remark}
The repository already contains most arithmetic pieces of this proof, but the
complete generic generated reflection library is one of the central Lean
milestones.  The report claims the mathematical schema, not an already
compiled implementation.
\end{remark}

\subsection{Bounded least selection}

Let \(v(q,c)\) be a primitive-recursive bit.  The checked bounded-search module
finds the least index below a bound at which a chosen bit convention succeeds.
A small adapter converts between ``zero means success'' and ``one means
success''.

\begin{theorem}[Canonical bounded selection]
\label{thm:least}
There is a primitive-recursive function \(\mu_v(q,B)\) such that \(\ipra\)
derives
\[
 \bigl(\exists c<B\;v(q,c)=1\bigr)
 \longrightarrow
 \left(
 \begin{array}{l}
 \mu_v(q,B)<B,\\[-2pt]
 v(q,\mu_v(q,B))=1,\\[-2pt]
 \forall j<\mu_v(q,B)\;v(q,j)=0
 \end{array}
 \right).
\]
\end{theorem}

\begin{proof}
Apply the checked least-zero search to the complement bit
\(1\mathbin{\dot-}v\).  Its specification supplies an accepted index no later
than any supplied accepted witness and rules out all earlier successes.  The
remaining steps are Boolean arithmetic and the existing order lemmas.
\end{proof}

\begin{principle}[Use search only at an escape boundary]
If a certificate is used only in a proof branch, apply existential elimination.
Invoke \cref{thm:least} only when a downstream term or public theorem must name
a code.  This avoids unnecessary normalization burdens and prevents a local
existence proof from being mistaken for global choice.
\end{principle}

\subsection{Finite master certificates}

Suppose that for each \(i<N\) a bounded decidable verifier \(v(i,c)\) has a
witness below \(B(i)\).  A common informal step says ``choose all witnesses''.
In intuitionistic arithmetic this should be made explicit.

\begin{theorem}[Master-certificate theorem]
\label{thm:master}
Assume \(\ipra\) derives
\[
       \forall i<N\;\exists c<B(i)\;v(i,c)=1.
\]
Then there is a primitive-recursive sequence code \(C(N)\) for which \(\ipra\)
derives
\[
       \forall i<N\;
       \bigl(\seq(C(N),i)<B(i)\wedge
             v(i,\seq(C(N),i))=1\bigr).
\]
The construction uses no choice principle.
\end{theorem}

\begin{proof}
Define \(m(i)=\mu_{v(i,\cdot)}(B(i))\) by bounded least search.  The premise and
\cref{thm:least} prove the validity of \(m(i)\) for arbitrary \(i<N\).  Define
\(C(0)\) as the empty sequence and append \(m(k)\) at stage \(k+1\).  The
sequence equations are primitive recursive.  Quantifier-free induction on
\(k\) proves the decoding equations for entries already appended.  For an
arbitrary \(i<N\), instantiate the decoding lemma and the validity theorem for
\(m(i)\).  Mixed-radix coding gives an alternative when all component bounds
are available in advance.
\end{proof}

The theorem resolves hidden finite choice in proof graphs, finite covers,
strategy tables, and compatibility diagrams.  Notice that decidability and a
bound are essential; an unbounded family of arbitrary existential statements
would be stronger.

\subsection{Finite monotone demand closure}

Adaptive proofs can have cyclic precision requirements: node \(i\) asks for a
resource determined by the resources requested by other nodes.  Unbounded
fixed-point language is unnecessary when each coordinate has a cap.

\begin{theorem}[Least finite demand closure]
\label{thm:demand-closure}
Let
\[
      P=\prod_{i<r}\{0,1,\ldots,C_i\}
\]
with coordinatewise order, and let \(F:P\to P\) be monotone.  Define
\[
 d^0=0,\qquad d^{n+1}=d^n\vee F(d^n),
\]
where \(\vee\) is coordinatewise maximum.  Put \(K=\sum_{i<r}C_i\).  Then:

\begin{enumerate}
\item \(d^K=d^{K+1}\), so \(d^K\) is a prefixed point of \(F\);
\item if \(e\in P\) satisfies \(F(e)\le e\), then \(d^K\le e\);
\item consequently \(d^K\) is the least prefixed point of \(F\).
\end{enumerate}
\end{theorem}

\begin{proof}
The sequence is increasing.  If \(d^{n+1}\ne d^n\), at least one coordinate
increases by at least one, so the potential
\(\Phi(d)=\sum_{i<r}d_i\) strictly increases.  The potential starts at zero
and is at most \(K\).  Hence a strict increase cannot occur at all of the first
\(K+1\) stages, and monotonicity of the update makes the first equality
permanent.  This proves stabilization by stage \(K\).

For minimality, prove by induction on \(n\) that \(d^n\le e\).  It is true for
zero.  If \(d^n\le e\), monotonicity gives \(F(d^n)\le F(e)\le e\), and hence
\(d^{n+1}=d^n\vee F(d^n)\le e\).  In particular \(d^K\le e\).

Every operation is finite and primitive recursive.  The formal proof uses
encoded arrays, bounded sums, order, and quantifier-free induction on the
iteration counter.  The strict-increase argument can be expressed by the
potential equation and the checked order library.
\end{proof}

\subsection{Finite strategy certificates}

A fixed bounded quantifier alternation can be represented without classical
choice.

\begin{definition}[Strategy certificate]
For a fixed formula
\[
 Q_0x_0<B_0\;Q_1x_1<B_1(\vec x_{<1})\cdots
 Q_sx_s<B_s(\vec x_{<s})\;A(\vec x),
\]
where \(A\) is reflected and each \(Q_i\) is \(\forall\) or \(\exists\), a
strategy certificate contains, at every existential node, a response as a
function of the preceding universal moves, and enough finite table data to
verify every leaf.
\end{definition}

\begin{theorem}[Bounded strategy realization]
\label{thm:strategy}
For each fixed bounded quantifier pattern there is a primitive-recursive
strategy verifier and a generated \(\ipra\) derivation from ``the strategy
verifies'' to the bounded formula.
\end{theorem}

\begin{proof}
Induct externally on the fixed quantifier pattern.  At an existential node,
decode the response selected by the strategy, verify its bound, and apply
existential introduction.  At a universal node, introduce an arbitrary bounded
move and select the corresponding certified subtree.  At a leaf, apply
\cref{thm:boolean-reflection}.  The only induction internal to \(\ipra\) checks
finite table lookup and bounds, so its invariant is quantifier-free.  No law of
excluded middle is used beyond reflected decidability of the bounded atoms.
\end{proof}

The theorem covers finite games, bounded minimax arguments, finite adversarial
verification, and many bounded classical case analyses.  It does not turn an
unbounded classical proof into a witness-producing intuitionistic proof.

\section{The FIBER compiler and its main theorem}
\label{sec:compiler}

\subsection{A small proof-producing intermediate language}

The source language is not intended to replace Lean or the mathematical prose
of a packet.  It records only the finite-use structure needed for compilation.
A program is a finite directed graph whose nodes have one of the following
forms.

\begin{longtable}{p{0.18\linewidth}p{0.72\linewidth}}
\toprule
Node & Meaning\\
\midrule
\endfirsthead
\toprule
Node & Meaning\\
\midrule
\endhead
\textsf{primitive} & Apply a fixed primitive-recursive operation and a local checked contract.\\
\textsf{pair}/\textsf{project} & Form or read a finite tuple or encoded sequence.\\
\textsf{map}/\textsf{fold} & Execute a bounded loop with a quantifier-free state invariant.\\
\textsf{search} & Use checked bounded least search for a reflected verifier.\\
\textsf{case} & Branch on a reflected Boolean atom and join the two derivations by disjunction elimination.\\
\textsf{exists} & Open a local certificate and seal it by existential elimination.\\
\textsf{strategy} & Verify and realize a fixed bounded quantifier game.\\
\textsf{observe} & Export an equivalence-invariant or metric-stable observation.\\
\textsf{residual} & Add an exact producer formula to the local context.\\
\textsf{call} & Invoke another compiled interface at a finite list of demanded queries.\\
\bottomrule
\end{longtable}

Each node declares its input and output arities, query transformer, certificate
bound, verifier bit, formula template, and local proof theorem.  Acyclic graphs
are topologically scheduled.  A graph with monotone cyclic demand first invokes
\cref{thm:demand-closure}; the resulting finite demand is then unrolled into an
acyclic schedule.

\subsection{The compiler invariant}

For every scheduled node \(v\), the compiler maintains:

\begin{enumerate}
\item a term tuple \(s_v\) containing the node state;
\item a quantifier-free state formula \(I_v(s_v)\);
\item a finite list \(\Gamma_v\) of residual and input formulas;
\item a derivation \(\Gamma_v\vdash_{\PRAAx}I_v(s_v)\);
\item a map from local certificate variables to their lexical scopes;
\item an observation relation guaranteeing that no noninvariant local code
      escapes.
\end{enumerate}

Bounded logical content is reflected into bits before entering a loop.
Quantifiers may occur in the external contract of a node, but the formula on
which the \(\ipra\) induction rule acts is always quantifier-free.

\subsection{Compilation theorem}

\begin{theorem}[FIBER compilation]
\label{thm:compilation}
Let \(P\) be a fixed well-formed FIBER program from interfaces
\(X_1,\ldots,X_m\) to \(Y\).  Suppose:

\begin{enumerate}
\item every primitive node is supplied with a local derivation in
      \(\Prf\;\PRAAx\);
\item every bounded formula used by a checker has a reflection derivation;
\item every loop carries a quantifier-free invariant with base and step
      derivations;
\item every cyclic demand operator is monotone and has explicit finite caps;
\item every exported operation respects the declared observation relation;
\item every external requirement is named by a residual formula.
\end{enumerate}

Then an effective metalevel compiler constructs:

\begin{enumerate}
\item fixed primitive-recursive symbols for demand, master-certificate bounds,
      verification, transformation, and observation;
\item a finite list of residual formulas \(R_P(p,q)\);
\item a Lean proof term of the literal type
\[
 R_P(p,q),\;\mathsf{Inputs}_P(p,q)
 \;\vdash_{\PRAAx}\;
 \mathsf{Output}_Y(p,q);
\]
\item and, when the output is certificate-independent, a proof in which all
      local and master certificates have been existentially eliminated.
\end{enumerate}

If derivations of all formulas in \(R_P(p,q)\) are later supplied, the closed
or client-context conclusion follows by \texttt{Prf.trans}.
\end{theorem}

\begin{proof}
The proof is by external induction over the finite schedule.

At a primitive node, instantiate its closed local theorem at the actual input
terms using \texttt{LeanProofs.IPRA.inst}, then weaken to the current context.
At a pairing or projection node, use the checked pairing equations.  At a
bounded map or fold, apply quantifier-free induction to the declared state
formula; the base and step derivations are part of the node contract.  At a
search node, instantiate \cref{thm:least}.  At a case node, reflect the bit,
derive the finite decidable disjunction, compile each branch, and join them by
\texttt{Prf.orE}.  At an existential node, compile the body under the shifted
certificate hypothesis and apply \texttt{Prf.exE}.  At a strategy node, use
\cref{thm:strategy}.  At an observation node, apply the operation congruence
or metric stability theorem and then seal the local certificate.

A call node first computes the finite demanded queries.  For each demand it
instantiates the called interface theorem.  The input derivations are then
substituted into the caller's local theorem using \texttt{Prf.trans}.  If a
single code is required, \cref{thm:master} forms the master certificate; if
not, the certificates are opened one at a time and sealed.

For monotone cyclic demand, \cref{thm:demand-closure} computes the least finite
closure.  The graph is unrolled for the resulting bounded number of instances,
so the remaining schedule is finite and acyclic.

The induction maintains the lexical-scope invariant.  A certificate variable
can disappear only through existential elimination; an exported term may
mention it only if the node is marked canonical and least search has produced
a closed term.  Therefore the final conclusion contains no accidental witness
escape.  Every unresolved producer encountered by the schedule is inserted as
a hypothesis, yielding exactly \(R_P(p,q)\).  Finally simultaneous cut
combines all node derivations and, when available, discharges residuals.
\end{proof}

\subsection{Composition}

\begin{theorem}[Compositionality]
\label{thm:composition}
If \(R:X\to Y\) and \(S:Y\to Z\) are FIBER realizers, then there is a FIBER
realizer \(S\circ R:X\to Z\).  Its demand is obtained by flattening the
\(R\)-demands for each \(Y\)-query demanded by \(S\); its residual list is the
union of the specialized residuals of \(R\) and \(S\); and its soundness proof
is obtained by substitution and simultaneous cut.
\end{theorem}

\begin{proof}
For a target query \(q\), compute \(\dem_S(q)=\langle y_0,\ldots,y_{k-1}\rangle\).
For each \(y_i\), compute \(\dem_R(y_i)\), concatenate the lists, and remove
no duplicates unless a proved normalization theorem is available.  A master
certificate supplies the required \(X\)-certificates.  The transformer first
runs \(T_R\) at each \(y_i\), bundles the resulting \(Y\)-certificates, and
then runs \(T_S\).  Bounds and all operations are primitive recursive.

In the object theory, instantiate the soundness theorem for \(R\) at every
\(y_i\).  Use these derivations to discharge the \(Y\)-validity hypotheses of
the soundness theorem for \(S\) by \texttt{Prf.trans}.  Extensionality or
metric stability composes by transitivity or the triangle inequality.  Local
certificates are sealed.  The remaining hypotheses are precisely the
specialized residuals.  Sequence concatenation may make associativity hold
only up to a provable coding equality, but the exported observation is
unchanged; normalizing all flattened demands to one canonical sequence code
makes composition strictly associative at the metadata level.
\end{proof}

\begin{corollary}[Reuse under arbitrary terms and contexts]
A compiled closed interface theorem can be specialized at arbitrary object
terms and used under arbitrary client hypotheses.
\end{corollary}

\begin{proof}
Apply \texttt{LeanProofs.IPRA.inst} to substitute the terms and then
\texttt{Prf.ofNil} or weakening to add the client context.  This is why reusable
lemmas should be stated either at arbitrary terms or in an explicitly
universally closed form.
\end{proof}

\subsection{Soundness versus completeness}

The compilation theorem is a soundness theorem: every accepted program yields
an \(\ipra\) derivation from its explicit residuals.  It does not assert that
every \(\ipra\) theorem has a convenient FIBER presentation, nor that every
classical proof can be mined automatically.  This asymmetry is intentional.
The source language is a disciplined fragment optimized for finite,
quantitative mathematics and auditability.

\subsection{Why the framework is not merely a checker generator}

A checker alone proves only a conditional statement of the form ``if this
finite code verifies, then the local conclusion holds.''  The compiler also
tracks:

\begin{itemize}
\item where the certificate comes from;
\item whether its existence is local, canonical, or residual;
\item how a downstream query changes upstream precision demand;
\item whether two independently produced certificates are compatible;
\item which observations are stable under representation changes;
\item which formulas are actually used as quantifier-free induction
      invariants;
\item and how all local proofs are assembled into the final object sequent.
\end{itemize}

These are exactly the points at which the existing packet outlines most often
stop.

\section{Residual producer calculus}
\label{sec:residuals}

A residual is not an informal promise that ``some effective constant exists.''
It is a fixed formula in the final object-language context.  This section gives
a normal form and a hierarchy of eliminations.

\subsection{Residual normal form}

A producer residual is written
\[
 \mathsf{Produce}_j(p,q)\;:\equiv\;
 \exists z<B_j(p,q)\;W_j(p,q,z)=1
\]
when a finite bound and verifier are known.  When no bound is known, the exact
formula records what is missing, for example:

\begin{align*}
&\exists H\;\forall P\;(
       \mathsf{Solution}(P)\to\mathsf{Height}(P)\le H),\\
&\exists N\;\forall n\ge N\;\mathsf{Error}(n)\le 2^{-k},\\
&\exists \eta>0\;\mathsf{Sep}(p,\eta),\\
&\exists L\;\mathsf{CompleteRepresentativeList}(p,L),\\
&\exists C,N_0\;\forall n\ge N_0\;\mathsf{Remainder}(n)\le C\rho(n).
\end{align*}

The residual formula must expose all parameters and conventions: strict versus
nonstrict bounds, zero-divisor behavior, representation of height, requested
precision, and whether the output is a list, a bound, or a named witness.

\subsection{Classes of residuals}

\begin{longtable}{p{0.18\linewidth}p{0.42\linewidth}p{0.31\linewidth}}
\toprule
Class & Available data & FIBER treatment\\
\midrule
\endfirsthead
\toprule
Class & Available data & FIBER treatment\\
\midrule
\endhead
R0: structural & Only substitution, weakening, conjunction, or finite cut is missing. & Eliminate directly with checked \texttt{Prf} combinators.\\
R1: bounded decidable & A bound \(B\), verifier bit \(W\), and object proof of existence are available. & Use checked least search; produce a canonical witness if needed.\\
R2: finite family & For each \(i<N\), an R1 producer is available. & Form one master certificate by \cref{thm:master}.\\
R3: monotone demand & A finite product of caps and a monotone demand operator are available. & Compute the least closure by \cref{thm:demand-closure}.\\
R4: metric coherence & Adjacent error bounds and a primitive-recursive summable tail majorant are available. & Derive a global Cauchy modulus by finite summation; see \cref{sec:analysis}.\\
R5: compactness substitute & A finite net, coverage checker, and continuity modulus are available. & Seal the local covering point and lift the finite calculation.\\
R6: bounded classical & A fixed bounded game or finite decidable case split is available. & Verify a finite strategy certificate and realize it intuitionistically.\\
R7: guarded selector & A positive separation or apartness margin is available. & Define a stable exact label or partial operation under the guard.\\
R8: genuinely external & No effective bound, modulus, representative theorem, or finite presentation is known. & Leave the exact producer formula in the final sequent.\\
\bottomrule
\end{longtable}

The hierarchy is not a scale of mathematical importance.  It separates
engineering omissions from genuine strength.  R0--R7 often disappear once the
correct finite formulation is chosen.  R8 is the honest boundary of the
application.

\subsection{Eliminating bounded residuals}

\begin{proposition}[Bounded producer elimination]
Suppose a compiled theorem has endpoint
\[
 \mathsf{Produce}(p,q),\Gamma\vdash_{\PRAAx}\Phi(p,q)
\]
with
\(\mathsf{Produce}(p,q)\equiv\exists z<B(p,q)\,W(p,q,z)=1\), and suppose
\(\Gamma\vdash_{\PRAAx}\mathsf{Produce}(p,q)\).  Then
\(\Gamma\vdash_{\PRAAx}\Phi(p,q)\).  If \(\Phi\) names the witness, replace
existential elimination by the least selector from \cref{thm:least}.
\end{proposition}

\begin{proof}
For a witness-independent conclusion, apply simultaneous cut or existential
elimination.  For a named witness, derive validity of the least accepted code
and substitute that term in the endpoint.  Both constructions are checked
object derivations.
\end{proof}

\subsection{From finite totality to a producer program}

A common packet proves, for arbitrary \(i<N\), that some local certificate
exists, but stops before constructing a uniform certificate array.  The
master-certificate theorem closes this gap when the verifier is bounded and
decidable.  In source-language terms the residual
\[
   \forall i<N\,\exists c<B(i)\;V(i,c)=1
\]
is compiled to the primitive-recursive producer
\[
   C(i)=\mu c<B(i)\,[V(i,c)=1],
\]
followed by sequence recursion.  This is a valid finite choice principle
because the least search is definable and its correctness is proved.  It
should not be generalized to unbounded or nondecidable relations.

\subsection{Residual factoring}

Residuals should be factored at the lowest reusable level.  For example, a
``root branch theorem'' may appear to require one monolithic residual, but it
usually decomposes into:

\begin{enumerate}
\item exact coefficient presentation;
\item a rational isolating box;
\item a derivative or resultant separation margin;
\item a polynomial interval-evaluation certificate;
\item and a requested precision.
\end{enumerate}

The last three may be bounded decidable data once the first two are supplied.
Factoring lets a later theorem discharge only the genuinely new producer and
reuse the finite checker.  The same principle applies to analytic constants,
finite covers, representatives of descent classes, and topological carriers.

\subsection{Non-eliminable examples}

The framework does not manufacture missing effective mathematics.  In
particular:

\begin{itemize}
\item Finiteness of a Diophantine set does not supply a computable height
      cutoff.  A complete bounded enumeration requires an effective height
      theorem or a complete list theorem.
\item Finiteness of a Selmer set does not supply rational points representing
      all locally soluble covering classes.  Representative production and
      compatibility remain separate formulas.
\item Compactness in an abstract completed space does not supply a finite net
      in the chosen coding.  A net and a coverage modulus must be produced.
\item Metastability for every control function does not by itself give a
      Cauchy rate.  A tail modulus or summable adjacent-motion estimate is new
      data.
\item A classical proof of \(\neg\neg\exists x A(x)\) does not yield a witness
      unless the conclusion is stable, bounded decidable, or accompanied by a
      producer.
\end{itemize}

In each case FIBER still contributes something useful: it turns the remainder
of the proof into a verified finite consumer and states the missing theorem in
a form that can be attacked independently.

\subsection{A theorem-status discipline}

Each theorem in an implementation should receive one of the following labels.

\begin{description}
\item[Checked object theorem.] A compiled \texttt{Prf PRAAx} declaration.
\item[Inherited checked theorem.] A named declaration from ConciseIPRA,
      IPRADerivations, or IPRABridge, with a pinned import path.
\item[Human-complete schema.] A full mathematical proof whose generic Lean
      implementation remains to be completed.
\item[Executable evidence.] Finite tests or exact computations; never a
      substitute for an object derivation.
\item[Conditional theorem.] A compiled or human proof with displayed residual
      producer formulas.
\item[Obstructed formulation.] A stronger requested interface is false or
      impossible; a corrected relational, metric, or guarded statement is
      supplied.
\item[Open.] Neither proof nor adequate residual reduction is available.
\end{description}

This vocabulary is used in the theorem ledger shipped with the archive.

\section{Effective analysis: from local certificates to coherent names}
\label{sec:analysis}

Many analysis packets correctly replace limits by rational approximants and
error bounds, but a repeated gap remains: separately certified approximants
are treated as one coherent real name.  FIBER distinguishes three levels:
pointwise approximation, finite-window stability, and a global Cauchy modulus.
Only the third is a completed name.

\subsection{Dyadic names and adjacent motion}

Represent a dyadic rational at precision \(n\) by a signed integer
\(a_n\), interpreted as \(a_n/2^n\).  Signed integers can be represented by the
existing difference terms in \texttt{IPRADerivations.Core.Integers}.  Define
the scaled discrepancy, for \(m\ge n\), by
\[
       \Delta(n,m)=\left|2^{m-n}a_n-a_m\right|.
\]
This is an integer quantity and therefore avoids division in the object
language.

\begin{theorem}[Dyadic adjacent-motion theorem]
\label{thm:dyadic}
Suppose a sequence of signed integers satisfies
\[
        |2a_k-a_{k+1}|\le 1
        \qquad\text{for every }k.
\]
Then for all \(m\ge n\),
\[
        \Delta(n,m)\le 2^{m-n}-1.
\]
Consequently
\[
        \left|\frac{a_n}{2^n}-\frac{a_m}{2^m}\right|<2^{-n}.
\]
Thus \(n\mapsto a_n/2^n\) is a Cauchy name with modulus \(n\mapsto n\).
\end{theorem}

\begin{proof}
Write \(r=m-n\) and induct on \(r\).  For \(r=0\), the discrepancy is zero,
which equals \(2^0-1\).  For the successor step, use
\begin{align*}
\Delta(n,m+1)
 &=|2^{m+1-n}a_n-a_{m+1}|\\
 &=|2(2^{m-n}a_n-a_m)+(2a_m-a_{m+1})|\\
 &\le 2\Delta(n,m)+|2a_m-a_{m+1}|\\
 &\le 2(2^{m-n}-1)+1\\
 &=2^{m+1-n}-1.
\end{align*}
The induction invariant is an inequality between integer-coded terms, hence
quantifier-free.  The object proof uses integer triangle inequality, powers of
two, order monotonicity, and quantifier-free induction.

Dividing the scaled inequality by \(2^m\) is only an explanatory
interpretation.  The formal endpoint can remain
\(\Delta(n,m)<2^{m-n}\), an integer inequality.  This is exactly the usual
Cauchy estimate after clearing denominators.
\end{proof}

\begin{corollary}[Independent certificates become coherent]
Suppose that for each precision \(n\) a local certificate produces \(a_n\), and
a pair certificate proves \(|2a_n-a_{n+1}|\le 1\).  Then the certificates may
be sealed while the sequence of observations exports a Cauchy name.
\end{corollary}

\begin{proof}
Use local existential elimination to obtain each adjacent pair only inside the
proof of its inequality.  The public conclusion depends solely on the
observations \(a_n,a_{n+1}\).  Apply \cref{thm:dyadic}.  No single global
certificate choosing all pointwise witnesses is necessary unless the name
itself must be returned as a code; in that case bounded least selection and a
master certificate provide one.
\end{proof}

\subsection{General summable-tail theorem}

The dyadic estimate is a special case of a finite telescoping principle.
Let \(e_k\ge0\) be integer or rational error majorants, and let \(T(n)\) be a
primitive-recursive tail certificate satisfying, for every \(m\ge n\),
\[
        \sum_{k=n}^{m-1}e_k\le T(n).
\]
If \(|x_{k+1}-x_k|\le e_k\), then
\[
        |x_m-x_n|\le T(n).
\]
The proof is finite induction using the triangle inequality and the bounded
sum equations.  A global rate therefore follows from an explicit summable
tail majorant; it does not follow from pointwise convergence or metastability
alone.

\subsection{Metastability is a separate interface}

A metastable contract has the form
\[
 \forall k,g\;\exists N\le\Phi(k,g)\;
 \forall i,j\in[N,N+g(N)]\;d(x_i,x_j)\le2^{-k}.
\]
For fixed \(k,g\) this is a finite query and fits FIBER directly: a certificate
contains \(N\) and a finite table or checker for the window.  It is often
provable when no useful global rate is known.  The framework keeps it as a
metastable interface.  Converting it to a Cauchy name requires additional
data, such as a summable adjacent-motion estimate or an independently supplied
tail modulus.  This prevents an invalid strengthening found in several
informal proof plans.

\subsection{Finite-net lifting}

Abstract compactness should be replaced by the finite data actually consumed
at a target precision.

\begin{theorem}[Finite-net lifting]
\label{thm:net}
Fix encoded domains \(X,Y\), discrepancy functions \(d_X,d_Y\), a function
\(f:X\to Y\), and a predicate \(P(y)\).  For a target query suppose there are:

\begin{enumerate}
\item a finite list \(z_0,\ldots,z_{L-1}\) and a coverage certificate
\[
       \forall x\in X\;\exists i<L\;d_X(x,z_i)\le\epsilon;
\]
\item a modulus certificate
\[
       d_X(x,z)\le\epsilon\longrightarrow
       d_Y(f(x),f(z))\le\delta;
\]
\item a stability theorem
\[
       P(f(z))\wedge d_Y(f(x),f(z))\le\delta
       \longrightarrow P_\delta(f(x));
\]
\item finite proofs of \(P(f(z_i))\) for all \(i<L\).
\end{enumerate}

Then \(\ipra\) proves \(P_\delta(f(x))\) for every encoded \(x\in X\).
The chosen net point is local and need not be exported.
\end{theorem}

\begin{proof}
Introduce an arbitrary \(x\) satisfying the domain predicate.  Open the
coverage existential to obtain an index \(i<L\) with the distance bound.
Instantiate the finite net theorem at \(i\), apply the modulus, then apply
stability.  The conclusion does not mention \(i\), so existential elimination
seals it.  The finite family of net-point proofs can be represented by a
master certificate or by a reflected bounded universal.  Every induction used
to verify list lookup or the finite family is quantifier-free after
Booleanization.
\end{proof}

\begin{remark}
This theorem does not assert that an arbitrary compact metric space has an
effective net in the selected coding.  Net production and coverage are
explicit hypotheses or residuals.  Once supplied, the remainder of the usual
compactness argument is a finite object-theory proof.
\end{remark}

\subsection{Guarded partial operations}

A statement such as \(d\ne0\) is generally too weak for stable numerical
division.  The useful interface is a positive apartness margin:
\[
             |d|\ge\eta>0.
\]
At target precision \(n\), the demand transformer requests enough precision on
\(d\) and the numerator to keep interval error below a fixed fraction of
\(\eta\).  Interval or Newton certificates then verify the reciprocal and the
quotient.  The same pattern applies to pivots, logarithms, square roots away
from zero, normalized eigenvectors, and branch labels.  The guard is a
residual if no effective lower bound is known.

\subsection{Modulus calculus}

FIBER realizers carry a demand map, which functions as a compositional modulus.
For metric interfaces, common operations have the following finite rules.

\begin{center}
\begin{tabularx}{0.96\linewidth}{Y{0.29\linewidth}X}
\toprule
Operation & Sufficient upstream demand\\
\midrule
addition & enough precision that the two errors sum below the target\\
multiplication on a bounded box & input errors divided by a certified magnitude bound\\
composition \(g\circ f\) & first apply the modulus of \(g\), then that of \(f\)\\
finite sum of \(L\) terms & target precision plus a binary-length allowance for \(L\)\\
division & target precision plus an apartness margin for the denominator\\
root label & cluster precision plus a positive separation gap\\
finite integral or quadrature & mesh and local error sufficient for the requested total error\\
\bottomrule
\end{tabularx}
\end{center}

These are not semantic continuity assumptions hidden in a typeclass.  Each is
a finite arithmetic contract with an object derivation.  Composition of the
demand maps is the proof-mining content of the framework.

\subsection{Worked analytic pattern: certified bisection}

For a rational polynomial \(p\), an interval \([l_0,r_0]\), and a target
precision \(n\), a certificate contains the bisection trace
\((l_k,r_k)_{k\le n}\).  The verifier checks:
\[
 p(l_k)\le0\le p(r_k),\qquad
 r_k-l_k=2^{-k}(r_0-l_0),
\]
and verifies that \((l_{k+1},r_{k+1})\) is the half selected by the exact
rational sign of the midpoint.  These are finite polynomial and order
calculations.  A uniqueness claim additionally requires a guard: for example,
a certificate that
\[
 p(y)-p(x)=(y-x)Q(x,y)
 \quad\text{and}\quad Q(x,y)\ge\lambda>0
\]
throughout the interval.  Then two zeros satisfy
\((y-x)Q(x,y)=0\), and positive cancellation gives \(x=y\).  This avoids an
appeal to the classical mean-value theorem and fits the checked ring/order
library.

\section{Algebraic objects, branches, and stable observations}
\label{sec:algebra}

Several packet difficulties are caused not by weak arithmetic but by asking
for a discontinuous name.  Ordered roots at a collision, a normalized
eigenvector for a repeated eigenvalue, a rotation axis at the identity, and an
exact sign for an arbitrary Cauchy real are not stable observations.  FIBER
separates the finite relation from a guarded label.

\subsection{A finite-use obstruction to exact sign}

\begin{theorem}[No unguarded finite-use exact sign selector]
\label{thm:no-sign}
No deterministic procedure that, on each input Cauchy name, inspects only
finitely many approximants can return the exact trichotomy
\(-1,0,+1\) correctly for every real represented by an arbitrary rapidly
convergent rational name.
\end{theorem}

\begin{proof}
Run the procedure on the zero name.  During this computation it inspects a
finite set of indices; let \(N\) exceed all of them.  Construct a valid name
for the positive real \(2^{-N-3}\) whose first \(N\) queried approximants are
identical to those of the zero name and whose later approximants converge with
the required rate.  The procedure follows the same finite computation and
returns the same output on both names, so it is wrong on one of them.  The same
argument works for a negative perturbation.
\end{proof}

The theorem is a continuity argument about representations, not a statement
that signs can never be decided.  A positive separation witness
\(|x|\ge2^{-k}\) makes the sign decidable from sufficiently precise
approximants.  Therefore the correct interface is
\[
  \mathsf{Sign}(x,k,s):
  |x|\ge2^{-k}\longrightarrow s=\operatorname{sign}(x),
\]
or a relational output ``\(x\) lies in this interval'' when no gap is known.

\subsection{Root clusters instead of global labels}

Let \(P(T)\) be a monic polynomial of fixed degree \(d\).  The coefficient
vector is stable under perturbation; an ordered list of roots is not stable at
collisions.  FIBER therefore uses one of three outputs.

\begin{description}
\item[Symmetric observation.] Elementary symmetric functions, power sums,
      resultants, or the companion matrix.  These remain exact at collisions.
\item[Cluster observation.] A finite list of disjoint boxes with multiplicities
      and a coverage verifier.  The ordering of boxes is irrelevant.
\item[Guarded labels.] Under a positive separation residual between clusters,
      a canonical ordering and individual root names can be selected.
\end{description}

\begin{theorem}[Stable cluster transport]
Suppose two valid certificates for a polynomial at mesh \(\epsilon\) give
cluster families \(\mathcal C\) and \(\mathcal C'\).  If both certificates
prove complete coverage with multiplicity and every cluster has diameter at
most \(\epsilon\), then the bipartite overlap relation has a finite
multiplicity-preserving matching after clusters within distance
\(2\epsilon\) are merged.  Every symmetric or \(2\epsilon\)-stable observation
therefore agrees.
\end{theorem}

\begin{proof}
Form the finite incidence graph whose vertices are clusters from the two
certificates and whose edges record nonempty overlap of their inflated boxes.
A root counted by one complete certificate must be counted by the other, and
multiplicity is conserved on each connected component.  The component totals
therefore agree.  Merging each component gives a common coarsening with the
same multiplicities.  A stable observation changes by at most the prescribed
modulus under the diameter and inflation bounds.  The graph, component labels,
and multiplicity equations are finite objects; a certificate can contain a
spanning forest and the checker can verify every equation.
\end{proof}

This theorem is deliberately relational.  It does not construct a globally
ordered root list at a collision.

\subsection{Guarded exact labels}

\begin{proposition}[Gap-guarded labeling]
Suppose a cluster certificate gives \(d\) singleton clusters and a rational
gap \(\eta>0\) such that distinct clusters are separated by at least \(\eta\).
At any precision finer than \(\eta/4\), lexicographic ordering of the rational
boxes is independent of the valid certificate.  Hence individual labels are
certificate-independent and may be exported.
\end{proposition}

\begin{proof}
Two valid boxes for the same root overlap after inflation by their error
radius.  Boxes for different roots cannot cross or exchange order because the
total uncertainty is less than half the certified gap.  Thus the overlap
matching is unique and preserves lexicographic order.  The proof consists of
rational order inequalities and finite matching verification.
\end{proof}

The same pattern applies to eigenvalues.  Without a gap, export the
characteristic polynomial, spectral clusters, or projectors onto a cluster.
With a gap and a normalization guard, export a labeled eigenvector.  For
rotations, export a unit quaternion modulo sign or a projective axis relation;
only away from the identity should an exact axis label be requested.

\subsection{Finite algebra certificates}

A powerful branch-free representation of a zero-dimensional algebra is a
finite spanning certificate.  Let
\[
 A=R[x_1,\ldots,x_n]/(f_1,\ldots,f_s)
\]
and suppose a certificate gives polynomials
\(b_0=1,b_1,\ldots,b_{D-1}\) and, for every \(i,j\), coefficients
\(m_{i,j,k}\) and quotient polynomials witnessing
\[
  x_i b_j-\sum_{k<D}m_{i,j,k}b_k\in(f_1,\ldots,f_s).
\]
All equalities are finite polynomial identities and can be checked by ring
normalization with supplied quotient witnesses.

\begin{theorem}[Finite-algebra observation]
From such a certificate, multiplication by each \(x_i\) is represented by a
\(D\times D\) matrix.  Every polynomial observation is represented by a
matrix polynomial.  Traces, determinants, characteristic polynomials,
resultants, and Newton power sums are finite, functorial observations.  If the
certificate is transported under a coefficient substitution or quotient map,
these observations transport by the corresponding matrix identities.
\end{theorem}

\begin{proof}
Closure under multiplication by the generators implies closure under every
monomial and hence every polynomial, by induction on term construction.  The
supplied quotient identities show that the action descends modulo the ideal.
Matrix multiplication represents composition of multiplication operators.
The determinant, trace, and characteristic polynomial are finite polynomial
expressions in the matrix entries.  Substitution preserves all verified
polynomial identities, so it preserves the matrix representation and its
polynomial observations.  No basis independence is needed for the spanning
claim; if a basis or exact dimension is required, add a linear-independence
certificate.
\end{proof}

This representation addresses several Puiseux, elimination, and algebraic
geometry packets.  It gives stable finite observations before any branch is
selected.

\subsection{Exact versus approximate Puiseux output}

A uniform exact Puiseux branch selector from arbitrary Cauchy coefficient
names is obstructed for the same reason as exact sign: arbitrarily small
coefficient perturbations can split or permute branches.  Three corrected
interfaces are available.

\begin{enumerate}
\item With exact algebraic or splitting data, output a finite list of Puiseux
      expansions and a permutation-invariant certificate.
\item With approximate coefficients, output clusters of truncated expansions
      and residual bounds.
\item Under a discriminant or resultant separation guard, refine clusters to
      labeled branches.
\end{enumerate}

At each finite truncation the Newton polygon, coefficient equations, and
residual checks are finite.  What fails without additional data is not the
finite algebra; it is stable global naming.

\subsection{Descent and representative relations}

Arithmetic geometry often produces a finite class set but then informally
chooses rational representatives.  FIBER uses a relation
\(R(c,w)\) between a class code and a representative witness.  A finite
strategy or viability certificate verifies compatibility along the proof
graph.  Public conclusions that depend only on the class seal the witness.
Conclusions that require an actual rational point retain
\(\exists w\,R(c,w)\) as a residual unless a representative producer is
proved.  This prevents Selmer classes from being silently identified with
rational covering points.

\section{Application patterns and repaired proof obligations}
\label{sec:domains}

The framework is useful only if it turns reported obstacles into precise local
tasks.  This section records the normal form for the main domains represented
in the repository.

\subsection{Combinatorics and finite algebra}

Finite identities, recurrences, and inequalities are the easiest class.  A
proof normally needs:

\begin{enumerate}
\item a primitive-recursive symbol for the sequence or finite object;
\item defining equations at arbitrary object terms;
\item a quantifier-free induction invariant;
\item a ring, order, finite-sum, or bounded-search derivation;
\item an explicitly universally closed endpoint.
\end{enumerate}

The checked library already supplies most of this.  FIBER adds value when a
packet uses a large finite certificate or several independently generated
witnesses.  The certificate is encoded, its checker reflected, and local
witnesses sealed.  Examples include finite graph colorings, tiling boundary
checks, coefficient identities, truncated series, recurrence tables, and
finite matrix calculations.

\subsection{Real and complex analysis}

An analytic object is represented by a family of finite queries.  A typical
interface contains:

\begin{itemize}
\item a precision or derivative-order query;
\item rational or dyadic approximation data;
\item a residual, interval, or coefficient verifier;
\item a local modulus or tail bound;
\item an observation relation identifying independent valid certificates;
\item guarded operations for division, logarithms, branches, and inverses.
\end{itemize}

The compiler never assumes that a pointwise family is coherent.  Coherence is
proved by adjacent motion, a summable tail, nested intervals, or a viable path.
Finite-net lifting replaces compactness.  Contour and quadrature arguments are
finite polygonal or mesh calculations plus explicit local error bounds.
Analytic continuation is represented by a finite chain of overlapping discs,
with transition certificates and a spanning-tree compatibility proof.

For special functions, identities divide into three layers:

\begin{description}
\item[Exact finite algebra.] Coefficient recurrences, truncated products,
      differential equations after coefficient comparison, and polynomial
      transformations.  These should be the first object-theory milestones.
\item[Effective local analysis.] Convergence in a guarded domain, with a
      displayed majorant and tail modulus.
\item[Global continuation or asymptotics.] A finite path/sector certificate
      plus effective constants.  Missing constants remain residuals rather
      than being imported as opaque host data.
\end{description}

This layering gives a realistic route for the Gamma, Beta, Barnes G,
hypergeometric, Meijer G, Lambert W, and normal-distribution packets.  It also
prevents a successful finite-series check from being reported as a complete
object proof of the analytic function.

\subsection{Probability}

General measure spaces conceal both infinitary data and discontinuity.
FIBER begins with finite rational laws:

\[
   \{(x_i,w_i):i<L\},\qquad w_i\in\mathbb Q_{\ge0},
   \qquad \sum_{i<L}w_i=1.
\]

Expectation, moments, convolution, and distribution functions at rational
thresholds are finite sums.  Approximation of an external law requires a
quantization certificate, tail-moment bounds, and, for discontinuous events,
an anti-concentration or boundary-mass modulus.  A central limit theorem at a
fixed error and threshold is then a finite query: truncate tails, discretize,
verify a characteristic-function or smoothing estimate, and account for all
errors.  A general measure-theoretic statement is not silently inferred from
the finite law.

For medians and quantiles, exact labels are unstable when an interval of
medians exists.  Export the median interval or a finite relation; choose a
canonical endpoint only when the theorem actually needs one and its bounded
search contract is proved.

\subsection{Topology and geometry}

Completed spaces are replaced by finite carriers:

\begin{itemize}
\item simplicial or cubical complexes;
\item rational polygonal curves;
\item finite covers and nerves;
\item grid approximations with refinement maps;
\item crossing, orientation, or winding certificates;
\item local chart tuples with overlap checks.
\end{itemize}

A topology theorem at a fixed query becomes a finite combinatorial theorem plus
an approximation producer.  For example, a Jordan-curve query may request a
grid fine enough relative to a separation modulus; a flood-fill certificate
then proves the finite inside/outside claim.  Invariance under refinement is an
observation theorem.  The existence of the required separation or approximation
modulus is a distinct producer.

Finite gluing does not require a global sheaf object.  Given local data on a
finite cover and equality or error certificates on overlaps, choose a spanning
tree of the overlap graph.  Transport observations to a root along the tree.
Exact cocycle compatibility gives equality; metric edge errors add along the
path.  Remaining cycle edges are consistency checks.  This is the
``observational descent'' principle in a form directly compilable by FIBER.

\subsection{Diophantine and arithmetic geometry}

A bounded search can verify all solutions below a supplied height.  It cannot
prove that there are no larger solutions without an effective height or
complete-list theorem.  The correct interface has two parts:

\begin{enumerate}
\item a finite enumerator and verifier parameterized by a bound \(H\);
\item a residual producer asserting that every relevant solution has height
      at most \(H\), or that a supplied list is complete.
\end{enumerate}

Once the producer is available, least search and master-certificate formation
can construct canonical lists, counts, and finite invariants.  Without it, the
finite consumer is still reusable and the missing arithmetic theorem is
isolated exactly.

For descent methods, class codes and rational representatives are separated.
A finite relation records which representatives realize which class.  A
backward viability certificate chooses a compatible path through a finite
relation only when every node has a certified continuation.  This resolves
nonunique local witnesses; it does not prove that every abstract class has a
rational representative.

\subsection{Prime distribution and ineffective constants}

An asymptotic theorem used at a concrete input should expose:

\begin{itemize}
\item the threshold \(N_0\);
\item all constants in the error bound;
\item the range of moduli or parameters;
\item and the arithmetic checker for the final inequality.
\end{itemize}

If a classical proof establishes existence of constants but gives no
primitive-recursive producer or bound, those constants are residuals.  A
finite computation that finds plausible constants for a range is evidence for
an instance, not a uniform proof.  When a verified explicit estimate is
available, its constants become ordinary input terms and the downstream
argument is compiled by order and bounded-sum lemmas.

\subsection{Computability and proof coding}

The checked bridge and endpoint inventory already emphasize exact theories,
exact formulas, literal tables, and replay provenance.  FIBER applies the same
discipline to proof-coding reductions:

\begin{enumerate}
\item keep the exact outer theory in the theorem type;
\item separate semantic equivalence from an object-theory derivation;
\item treat a machine table, polynomial, or proof predicate as pinned data;
\item prove local arithmetic in the ordinary portable calculus;
\item translate at the Foundation or application boundary;
\item record directions, premises, and double-negation strength exactly.
\end{enumerate}

A semantic Lean theorem about a machine and an \(\ipra\) proof of a coded
nonhalting statement are different endpoints.  The framework does not merge
them in a status ledger.

\subsection{Classical arguments}

FIBER supports three constructive replacements for classical reasoning.

\begin{enumerate}
\item \textbf{Reflected bounded cases.}  A decidable finite atom yields an
      object proof of the finite disjunction; branches are joined by
      \texttt{orE}.
\item \textbf{Finite strategy certificates.}  Fixed bounded alternation is
      realized by \cref{thm:strategy}.
\item \textbf{Stable negative conclusions.}  A negative translation can be
      used when the target is stable and the translated assumptions are
      supplied; no witness extraction is claimed.
\end{enumerate}

An unbounded classical existence theorem that is neither stable nor equipped
with a producer remains outside these mechanisms.  This is an honest logical
boundary, not an implementation inconvenience.

\subsection{Representative repaired endpoints}

The following endpoints illustrate how a packet should be restated.

\begin{longtable}{p{0.24\linewidth}p{0.33\linewidth}p{0.34\linewidth}}
\toprule
Informal claim & Unsafe formulation & FIBER endpoint\\
\midrule
\endfirsthead
\toprule
Informal claim & Unsafe formulation & FIBER endpoint\\
\midrule
\endhead
A root has a branch & Return an exact global label from arbitrary Cauchy coefficients & Return a verified cluster; under a separation residual, refine to a label.\\
Approximants define a real & Prove each precision independently & Add adjacent-motion or nested-interval certificates and derive a Cauchy modulus.\\
Compactness gives a minimizer & Select a point from an abstract compact space & Supply a finite net, coverage and continuity moduli; return an approximate minimizer, with exactness only under a gap.\\
A finite set can be enumerated & Use finiteness to invent a height cutoff & Parameterize the verifier by \(H\); retain a complete-height/list residual.\\
A nonzero pivot can be divided by & Treat \(d\ne0\) as a numerical modulus & Require \(|d|\ge\eta>0\), or return a branch relation that avoids division.\\
A bounded proof graph has witnesses & Choose one witness at each node & Use least bounded search and one master certificate, or seal local witnesses.\\
A finite class has a representative & Identify class codes with rational points & Use a representative relation and retain representative totality as a producer.\\
A metastable sequence converges with rate & Infer a Cauchy modulus from finite windows & Export metastability; add a summable tail or explicit rate producer for a Cauchy name.\\
\bottomrule
\end{longtable}

\section{Lean formalization architecture}
\label{sec:lean}

\subsection{Target repository location}

A complete implementation should live under
\begin{center}
\texttt{Logic/PrimitiveRecursiveArithmetic/Derivations/Lean/IPRADerivations/FIBER/}
\end{center}
with a small public import
\texttt{IPRADerivations.FIBER}.  The compiler metadata and source generator may
live under \texttt{Tools/}, but every accepted theorem must end in the ordinary
\texttt{Prf PRAAx} calculus.  Foundation translations belong in a separate
bridge module so the portable library remains independent of application
closures.

\subsection{Existing declarations to reuse}

The implementation should import, rather than duplicate, the following checked
components.

\begin{longtable}{p{0.35\linewidth}p{0.56\linewidth}}
\toprule
Module/declaration & Use in FIBER\\
\midrule
\endfirsthead
\toprule
Module/declaration & Use in FIBER\\
\midrule
\endhead
\texttt{ConciseIPRA.Deduction} & \texttt{Prf.subst}, \texttt{weaken}, \texttt{trans}, \texttt{exE}, logical constructors, equality substitution.\\
\texttt{IPRADerivations.Core.Basic} & \texttt{inst}, open quantifier-free induction \texttt{ind}, arithmetic symbols and congruence.\\
\texttt{Core.Semiring}, \texttt{Core.Order}, \texttt{Core.Monus} & State arithmetic, bounds, comparison, monotonicity, cancellation.\\
\texttt{Core.Pairing} & Pairing, projections, and simultaneous recursion for sequence coding.\\
\texttt{Core.Search} & \texttt{Search.mu}, \texttt{search\_spec}, \texttt{search\_found}, and \texttt{search\_minimal}.\\
\texttt{Core.BoundedSums}, \texttt{Core.MultiSums} & Boolean folds, finite potentials, and master-verifier aggregation.\\
\texttt{Core.Integers}, \texttt{Core.Powers} & Dyadic names and the adjacent-motion theorem.\\
\texttt{Core.LinearAlgebra} & Matrix observations for finite algebras and spectral clusters.\\
\texttt{IPRABridge.OpenSequents}, \texttt{ScopedSequents} & Transport of completed open interface theorems to Foundation clients.\\
\bottomrule
\end{longtable}

\subsection{Module dependency order}

\begin{enumerate}
\item \texttt{FIBER/Kernel.lean}: structural aliases for residual discharge,
      specialization, local sealing, and context replacement.  This file
      should contain no arithmetic.
\item \texttt{FIBER/Bool.lean}: Boolean operations and reflection for atomic
      equality/order predicates.
\item \texttt{FIBER/Fold.lean}: bounded conjunction/disjunction folds and
      indexing lemmas.
\item \texttt{FIBER/BoundedFormula.lean}: external syntax for a restricted
      bounded formula language, generated checker terms, and reflection.
\item \texttt{FIBER/Sequence.lean}: one canonical bounded sequence coding,
      append, lookup, concatenation, and mixed-radix alternative.
\item \texttt{FIBER/SearchAdapter.lean}: one-success adapter on top of
      \texttt{Core.Search}.
\item \texttt{FIBER/Master.lean}: least local witnesses and the
      master-certificate theorem.
\item \texttt{FIBER/Demand.lean}: capped product orders, potential, monotone
      closure, and least prefixed point.
\item \texttt{FIBER/Strategy.lean}: fixed bounded games and strategy
      realization.
\item \texttt{FIBER/Interface.lean}: metadata structures at the Lean
      metalevel and corresponding object formula constructors.
\item \texttt{FIBER/Compile.lean}: proof-producing recursion over the finite
      source graph.
\item \texttt{FIBER/Analysis/Dyadic.lean} and
      \texttt{FiniteNet.lean}: the two principal cross-domain examples.
\item \texttt{FIBER/Algebra/Clusters.lean}: stable cluster observations and
      gap-guarded labels.
\item \texttt{FIBER/Checks/*.lean}: small closed examples, \texttt{\#print
      axioms}, and bridge round trips.
\end{enumerate}

This order deliberately completes the object kernel before a large metadata
language is introduced.

\subsection{Structural kernel signatures}

The first file can use the current API almost verbatim.

\begin{lstlisting}[language={},caption={Core structural targets}]
import IPRADerivations.Core.Basic

namespace LeanProofs.IPRA.FIBER

open LeanProofs.ConciseIPRA
open LeanProofs.IPRA
open Syn
open scoped Fm

variable {n : Nat}
variable {Gamma Delta : List (Fm PTm n)}
variable {phi : Fm PTm n}

/-- Simultaneous residual discharge. -/
theorem discharge
    (h : Gamma |- [PRAAx] phi)
    (hs : forall psi, psi in Gamma -> Delta |- [PRAAx] psi) :
    Delta |- [PRAAx] phi :=
  h.trans hs

/-- Reuse a closed theorem at arbitrary terms and hypotheses. -/
theorem specialize
    {k : Nat} {theta : Fm PTm k}
    (h : ([] : List (Fm PTm k)) |- [PRAAx] theta)
    (sigma : Fin k -> PTm n) :
    Gamma |- [PRAAx] theta.subst sigma :=
  LeanProofs.IPRA.inst h sigma

/-- Certificate sealing is the native existential-elimination rule. -/
theorem seal
    {A : Fm PTm (n + 1)} {B : Fm PTm n}
    (hex : Gamma |- [PRAAx] Fm.ex A)
    (hbody : A :: Gamma.map Fm.shift |- [PRAAx] B.shift) :
    Gamma |- [PRAAx] B :=
  Prf.exE hex hbody

end LeanProofs.IPRA.FIBER
\end{lstlisting}

The ASCII notation in the distributed Lean file uses the repository's actual
Unicode/scoped notation.  The proof terms are direct aliases, so this is the
first file to elaborate and audit.

\subsection{One-success search adapter}

\texttt{Core.Search} searches for a zero of a binary term \(f(i,x)\).  FIBER
verifiers conventionally accept when \(v(i,x)=1\).  Define
\[
       f_v(i,x)=1\mathbin{\dot-}v(i,x).
\]
Then instantiate \texttt{Search.search\_found} and prove:

\begin{lstlisting}[language={},caption={Target theorem for the search adapter}]
theorem leastAccepted_found
    {v : PTm 2} {N i x : PTm n}
    (hvBit : BitValued v)
    (hi : Gamma |- [PRAAx] lt i N)
    (hacc : Gamma |- [PRAAx] inst2 v i x == num 1) :
    Gamma |- [PRAAx]
      lt (leastAccepted v N x) N /\
      (inst2 v (leastAccepted v N x) x == num 1 /\
       le (leastAccepted v N x) i)
\end{lstlisting}

Here \texttt{BitValued} should be a proved object formula, not a host
proposition silently used to rewrite values.  The implementation needs the
existing monus and parity/Boolean lemmas to show that the complement is zero
exactly when \(v=1\).

\subsection{Bounded formula syntax}

A convenient external datatype is:

\begin{lstlisting}[language={},caption={Metalevel fixed bounded formulas}]
inductive BFm (n : Nat)
  | eq   : PTm n -> PTm n -> BFm n
  | lt   : PTm n -> PTm n -> BFm n
  | and  : BFm n -> BFm n -> BFm n
  | or   : BFm n -> BFm n -> BFm n
  | imp  : BFm n -> BFm n -> BFm n
  | bEx  : PTm n -> BFm (n+1) -> BFm n
  | bAll : PTm n -> BFm (n+1) -> BFm n
\end{lstlisting}

Define \texttt{BFm.toFm} and a compiler \texttt{BFm.bit : BFm n -> PTm n}.
The main theorem is an external recursion producing a proof:

\begin{lstlisting}[language={},caption={Reflection endpoint}]
def BFm.reflect (A : BFm n) :
  ([] : List (Fm PTm n)) |- [PRAAx]
    (A.bit == num 1 <-> A.toFm)
\end{lstlisting}

The bounded-quantifier cases should call fixed fold theorems rather than place
\texttt{A.toFm} in an induction predicate.  The generated proof is accepted
because \texttt{BFm} is a Lean metatype; there is no object-level universal
truth predicate.

\subsection{Interface metadata versus object theorems}

A Lean structure may organize the generator, but proof fields must be genuine
\texttt{Prf} values.

\begin{lstlisting}[language={},caption={Schematic interface record}]
structure Interface where
  arity      : Nat
  queryOK    : PR (arity + 1)
  certBound  : PR (arity + 1)
  verifier   : PR (arity + 2)
  observe    : PR (arity + 2)
  validFm    : Fm PTm (arity + 2)
  soundFm    : Fm PTm (arity + 2)
  sound      : ([] : List (Fm PTm (arity + 2)))
                 |- [PRAAx] soundFm
  -- plus equivalence/stability and substitution equations
\end{lstlisting}

Merely declaring a field named \texttt{sound} is not a proof of a concrete
instance.  Every instance must construct it without \texttt{sorry}, an
unproved \texttt{axiom}, or a host-level semantic cast.

\subsection{Demand closure target}

Use a fixed number of coordinates in the first implementation.  Encode a
vector by nested pairing or a bounded sequence.  Define:

\begin{lstlisting}[language={},caption={Demand closure theorem shape}]
structure DemandSystem (r : Nat) where
  cap      : Fin r -> PTm n
  step     : Fin r -> PTm (n + r)
  monotone : ObjectMonotoneProof step

theorem closure_fixed
    (S : DemandSystem r) :
    ([] : List (Fm PTm n)) |- [PRAAx]
      IsLeastPrefixedPoint S (iterate S (sumCaps S))
\end{lstlisting}

For a first checked milestone, take \(r=2\).  This exercises pairing,
coordinatewise maximum, the potential argument, and open induction without
requiring a generic variable-length array library.  Generalize only after the
fixed-arity proof is stable.

\subsection{Dyadic theorem target}

The formal endpoint should avoid rational division:

\begin{lstlisting}[language={},caption={Cleared-denominator Cauchy endpoint}]
/-- a : precision and parameter |-> signed numerator -/
theorem dyadic_scaled_cauchy
    (a : ZTerm 2)
    (hstep : AdjacentUnitProof a) :
    ([] : List (Fm PTm 3)) |- [PRAAx]
      (le (var 0) (var 1) -->
       zabs (zsub (zscalePow2 a (var 0) (var 1))
                  (a.at (var 1) (var 2)))
         < pow two (monus (var 1) (var 0)))
\end{lstlisting}

The implementation should first prove the exact geometric bound
\(2^{m-n}-1\), then derive the strict inequality.  This module provides a
small but meaningful demonstration that a framework theorem closes a reported
coherence gap in the object theory.

\subsection{Compiler implementation strategy}

There are two reasonable stages.

\paragraph{Stage A: trusted Lean generator.}
Represent a well-formed finite program as an inductive Lean datatype whose
constructors contain the local \texttt{Prf} proofs.  Define
\texttt{compile} by recursion, returning a \texttt{Prf} term.  Lean's kernel
checks the result; the generator itself need not be verified as an object
program.

\paragraph{Stage B: certificate-producing external frontend.}
A Python or Lean elaborator may parse JSON metadata, synthesize primitive
recursive terms, and emit Lean source.  Its output is still checked by Lean.
The external tool is not in the trusted base.  Keep generated formulas and
node IDs deterministic so proof artifacts can be diffed and replayed.

Do not begin with an internal self-interpreter.  The diagonal obstruction and
the size of syntax-substitution proofs make that interface both impossible in
the strongest form and unnecessary for the repository's goals.

\subsection{Required checks}

Each completed module should pass:

\begin{enumerate}
\item \texttt{lake build +IPRADerivations.FIBER.<Module>};
\item a no-placeholder scan for \texttt{sorry}, \texttt{admit}, and new
      theorem-specific axioms;
\item \texttt{\#print axioms} on every public endpoint;
\item at least one open-term specialization test;
\item at least one nonempty-context test;
\item a native-scalar transfer test via \texttt{pra\_iff\_native};
\item and, for a Foundation client, an open-sequent bridge round trip.
\end{enumerate}

Lean's standard metatheoretic axioms may appear in the report.  The current
checked repository already uses \texttt{propext},
\texttt{Classical.choice}, and \texttt{Quot.sound}.

The important audit is that no new object axiom or placeholder enters the
theorem.

\section{Claim ledger and remaining work}
\label{sec:status}

\subsection{What is established in this report}

The following results are supplied with complete human proofs or proof
reductions:

\begin{enumerate}
\item the fibered finite-query interface and realizer definitions;
\item local certificate sealing through existential elimination;
\item composition through substitution and simultaneous cut;
\item the fixed-formula Boolean reflection construction, with explicit
      quantifier-free induction accounting;
\item bounded least accepted-certificate selection from the checked search
      theorem;
\item master-certificate formation for finite bounded decidable families;
\item least finite monotone demand closure with a sharp potential bound;
\item realization of fixed bounded strategy certificates;
\item the FIBER compilation theorem by induction over the source schedule;
\item the residual-elimination hierarchy;
\item the dyadic adjacent-motion/Cauchy theorem;
\item finite-net lifting with local witness sealing;
\item guarded bisection and exact branch uniqueness from a positive factor
      bound;
\item the finite-use impossibility of unguarded exact sign selection;
\item the diagonal impossibility of a primitive-recursive universal evaluator;
\item stable cluster, finite-algebra, and representative-relation interfaces.
\end{enumerate}

The exact reference algorithms included in the archive test mixed-radix
bundling, least selection, demand closure, strategy trees, dyadic bounds,
finite-net lifting, and certificate escape behavior.

\subsection{What is inherited, not reproved}

The report relies on checked repository declarations for the object calculus,
substitution, weakening, simultaneous cut, quantifier-free induction,
arithmetic, order, search, pairing, bounded sums, powers, integers, and several
linear-algebra constructions.  It also relies on the checked bridge for
transport to Foundation clients.  These dependencies are named in the source
audit.  They were inspected through the GitHub snapshot but were not rebuilt
inside this artifact environment.

\subsection{What is not yet kernel-checked}

No new Lean theorem in this archive was compiled.  In particular, the
following remain implementation milestones:

\begin{enumerate}
\item a complete object proof of the bounded conjunction/disjunction fold
      reflection theorems;
\item the external \texttt{BFm.reflect} proof generator;
\item the one-success adapter over \texttt{Core.Search};
\item a canonical sequence coding suitable for generated master certificates;
\item the generic master-certificate theorem;
\item fixed-arity and then variable-arity demand closure;
\item the bounded strategy compiler;
\item the typed graph compiler returning literal \texttt{Prf PRAAx} terms;
\item the integer absolute-value and scaling lemmas needed for the dyadic
      theorem;
\item finite-net and cluster examples in the object theory;
\item integration builds and axiom audits.
\end{enumerate}

The distributed Lean files therefore use comments and theorem targets for
unfinished arithmetic.  They contain no claim that a source file without
\texttt{sorry} has been elaborated.

\subsection{Research boundaries not removed by the framework}

FIBER resolves representational and compositional obstacles.  It does not by
itself prove:

\begin{itemize}
\item general effective Faltings, Siegel, Roth, or Mordell--Weil generation;
\item effective Chebotarev, Siegel--Walfisz, or Bombieri--Vinogradov constants
      beyond what is explicitly supplied;
\item a finite presentation or modulus for every abstract compact space;
\item a Cauchy rate from metastability alone;
\item an exact branch, sign, eigenvector, or axis selector at degeneracy;
\item rational representatives for arbitrary abstract descent classes;
\item a universal primitive-recursive evaluator;
\item or automatic extraction from every classical proof.
\end{itemize}

For each boundary the framework gives either a corrected theorem or an exact
producer formula.

\subsection{Recommended first integration slice}

A realistic first pull request should be intentionally small:

\begin{enumerate}
\item add \texttt{FIBER/Kernel.lean} and compile its structural aliases;
\item add the one-success search adapter and a two-query master certificate;
\item prove a fixed two-coordinate demand closure theorem;
\item prove the dyadic adjacent-motion theorem for one explicit signed-integer
      encoding;
\item compile a miniature example that consumes two independently certified
      approximants and exports a Cauchy estimate;
\item run \texttt{\#print axioms}, native transfer, and an open-sequent bridge
      check.
\end{enumerate}

This slice exercises the distinctive contribution of the framework without
first requiring a generic formula compiler.  Once it is checked, the bounded
formula and source-graph generators can be built against a stable kernel.

\subsection{Acceptance criteria for a packet migrated to FIBER}

A packet should count as migrated only when:

\begin{enumerate}
\item its main target is stated as an explicit ordinary \(\ipra\) formula;
\item every primitive symbol and representation convention is fixed;
\item the finite query and demand map are explicit;
\item every checker has an object soundness/reflection derivation;
\item every induction predicate is displayed and quantifier-free;
\item all certificate scopes are classified and enforced;
\item all missing data appear as residual formulas;
\item the endpoint is a compiled \texttt{Prf PRAAx} declaration;
\item its axiom audit and build command are recorded;
\item and any Foundation or application endpoint is obtained through a checked
      bridge or an independently stated theorem.
\end{enumerate}

\section{Conclusion}

The ProveIt \(\ipra\) development has reached a point where the central
problem is architectural rather than foundational.  The ordinary calculus,
its small native core, a bridge to Foundation, and a substantial object
arithmetic library already exist.  The proof packets show that a large range
of mathematics can be reorganized into finite quantitative claims.  What is
needed is a reliable path from those claims to literal object derivations.

FIBER-iPRA supplies that path in a deliberately conservative form.  Queries
index finite certificate fibers.  Demand maps expose finite use.  Local
certificates are sealed.  Public codes are selected only by proved bounded
search.  Noncanonical objects export stable observations rather than
impossible labels.  Bounded loops use quantifier-free state invariants.
Composition is substitution and simultaneous cut, not semantic wishful
thinking.  Missing effective information is a formula, not an unstated host
parameter.

The framework also clarifies why several difficult-looking requests should be
reformulated.  A global exact branch at a collision, a Cauchy rate from mere
metastability, a height cutoff from abstract finiteness, and an internal
universal primitive-recursive evaluator are not reasonable implementation
milestones.  Cluster observations, summable-tail contracts, explicit height
producers, and metalevel specialization are the correct replacements.

The next decisive step is not another broad packet.  It is a small checked
kernel: search adaptation, finite certificate bundling, demand closure, and
the dyadic coherence theorem, all as \texttt{Prf PRAAx}.  Once these are in
place, the repository's many finite proof plans can share one audited
compilation route, and the remaining deep residuals will stand out as the
mathematics they genuinely are.


\appendix
\section{Formal source-language sketch}
\label{app:language}

This appendix gives a compact mathematical specification of the proposed
source language.  It is intentionally first-order and finite.

\subsection{Types}

Base types are \(\mathsf{Nat}\), \(\mathsf{Bit}\), and bounded sequences
\(\mathsf{Seq}(B,L)\).  Product types are encoded by pairing.  A metric
observation type adds a discrepancy program; a quotient observation type adds
an equivalence verifier and congruence proofs.  Types are metadata only: after
compilation every value is a natural-number term.

\subsection{Terms}

\[
\begin{array}{rcl}
t &::=& x \mid 0 \mid S(t) \mid f(t_1,\ldots,t_k)\\
  &\mid& \mathsf{pair}(t,u) \mid \mathsf{get}(t,i)\\
  &\mid& \mathsf{fold}(N,s_0,\lambda i\,s.\,h(i,s))\\
  &\mid& \mathsf{least}(N,\lambda i.\,b(i)).
\end{array}
\]
Every function symbol is a fixed \(\mathsf{PR}\) code generated externally or
imported from the checked library.

\subsection{Commands}

\[
\begin{array}{rcl}
P &::=& \mathsf{return}\;t\\
  &\mid& \mathsf{let}\;x=t\;\mathsf{in}\;P\\
  &\mid& \mathsf{call}\;R(q)\;\mathsf{as}\;x\;\mathsf{in}\;P\\
  &\mid& \mathsf{open}\;\exists x\,A(x)\;\mathsf{as}\;x\;\mathsf{in}\;P\\
  &\mid& \mathsf{if}\;b\;\mathsf{then}\;P\;\mathsf{else}\;P\\
  &\mid& \mathsf{loop}\;N\;\mathsf{state}\;s\;\mathsf{invariant}\;I\;\mathsf{do}\;P\\
  &\mid& \mathsf{require}\;R\;\mathsf{in}\;P\\
  &\mid& \mathsf{observe}\;o\;\mathsf{mod}\;E.
\end{array}
\]

A typing judgment has the form
\[
   \Sigma;\Lambda;\mathcal R\vdash P:(Y,q),
\]
where \(\Sigma\) is the term environment, \(\Lambda\) is the stack of local
certificate scopes, and \(\mathcal R\) is the residual list.  The key escape
rule is:
\[
 \frac{\Sigma,x;\Lambda,x;\mathcal R\vdash P:B
       \qquad x\notin\mathrm{FV}(B)}
      {\Sigma;\Lambda;\mathcal R\vdash
       \mathsf{open}\;\exists x A\;\mathsf{as}\;x\;\mathsf{in}\;P:B}.
\]
An observation may mention \(x\) only if an equivalence/stability theorem
proves independence, in which case the compiled conclusion is phrased through
the observation relation rather than code equality.

\subsection{Compilation judgment}

The compiler returns
\[
 \mathsf{compile}(P)=
 (\dem_P,B_P,V_P,T_P,O_P,\mathcal R_P,\pi_P),
\]
where \(\pi_P\) is a Lean value whose type is the appropriate
\(\Prf\;\PRAAx\) sequent.  The correctness theorem is therefore intrinsic:
there is no separate untrusted claim that the generated proof is valid.  Lean
accepts or rejects \(\pi_P\).

\subsection{Operational reference semantics}

A standard natural-number evaluator is useful for tests.  It maps a query and
certificate to either rejection or an observation.  The evaluator has no
logical authority.  The object soundness theorem is generated independently
from local derivations.  Keeping the executable semantics and proof compiler
separate makes it easier to detect a checker whose host behavior is correct but
whose object reflection theorem is missing.

\section{Lean theorem target catalogue}
\label{app:signatures}

The following catalogue is a dependency-ordered checklist.  Names are
provisional, but the conclusion type of every proof theorem should contain
\texttt{Prf PRAAx} explicitly.

\small
\begin{longtable}{Y{0.35\linewidth}Y{0.53\linewidth}}
\toprule
Target & Intended statement\\
\midrule
\endfirsthead
\toprule
Target & Intended statement\\
\midrule
\endhead
\texttt{Kernel.discharge} & Alias of \texttt{Prf.trans} specialized to ordinary \(\ipra\).\\
\texttt{Kernel.specialize} & Alias of \texttt{IPRA.inst}: closed open-formula theorem specialized at arbitrary terms and context.\\
\texttt{Kernel.seal} & Alias of \texttt{Prf.exE}.\\
\texttt{Bool.bit\_cases} & A bit-valued term is zero or one, with object derivation.\\
\texttt{Fold.any\_sound} & Accepted bounded disjunction fold yields an index and reflected body.\\
\texttt{Fold.any\_complete} & A supplied bounded witness forces the disjunction fold to one.\\
\texttt{Fold.all\_sound} & Conjunction fold one implies every bounded reflected body.\\
\texttt{Fold.all\_complete} & Bounded universal body implies conjunction fold one.\\
\texttt{BFm.reflect} & For fixed bounded syntax \(A\), \([A]=1\leftrightarrow A\).\\
\texttt{Search.leastAccepted\_found} & One-success wrapper over \texttt{Search.search\_found}.\\
\texttt{Search.leastAccepted\_minimal} & No earlier accepted index.\\
\texttt{Sequence.lookup\_append} & Appended entry and prior entries decode correctly.\\
\texttt{Master.valid} & Pointwise bounded totality yields one valid master code.\\
\texttt{Demand.iter\_mono} & Demand iteration is increasing.\\
\texttt{Demand.potential\_strict} & A nonfixed update strictly raises the coordinate sum.\\
\texttt{Demand.stabilizes} & Iteration stabilizes by the sum of caps.\\
\texttt{Demand.least} & Stabilized vector is below every prefixed point.\\
\texttt{Strategy.realize} & Verified fixed bounded strategy implies its quantified formula.\\
\texttt{Interface.localSound} & Valid demanded inputs and residuals imply valid output.\\
\texttt{Interface.obsIndependent} & Two valid output certificates give equivalent observations.\\
\texttt{Compile.sound} & Compiled program returns the final \texttt{Prf PRAAx} sequent.\\
\texttt{Compile.compose} & Composition of realizers and proof terms.\\
\texttt{Dyadic.scaledBound} & Adjacent unit discrepancies imply \(\Delta(n,m)\le2^{m-n}-1\).\\
\texttt{Dyadic.cauchy} & Cleared-denominator Cauchy inequality.\\
\texttt{FiniteNet.lift} & Coverage, modulus, stability, and finite checks imply the global approximate property.\\
\texttt{Clusters.commonCoarsening} & Two complete cluster certificates admit a verified common coarsening.\\
\texttt{Clusters.labelUnderGap} & Positive gap makes labels certificate-independent.\\
\texttt{Bridge.fiberOpen} & Transport a compiled open FIBER sequent to Foundation syntax.\\
\bottomrule
\end{longtable}
\normalsize

\subsection{Small regression examples}

Before a domain-scale application, the implementation should close these
examples:

\begin{enumerate}
\item choose the least even number in a supplied bounded interval;
\item combine two local divisibility witnesses into a certificate-independent
      gcd observation;
\item compute a two-coordinate cyclic demand closure;
\item verify a depth-two \(\forall\exists\) finite strategy;
\item form a master certificate for three independent polynomial evaluations;
\item export a dyadic Cauchy estimate from independently certified adjacent
      approximants;
\item lift an inequality from a rational finite net to all points represented
      by a coverage certificate;
\item transport one open theorem through IPRABridge and back.
\end{enumerate}

\section{Consolidated difficulty matrix}
\label{app:matrix}

\footnotesize
\setlength{\tabcolsep}{4pt}
\begin{longtable}{Y{0.17\linewidth}Y{0.24\linewidth}Y{0.27\linewidth}Y{0.15\linewidth}}
\toprule
Difficulty & Invalid shortcut & FIBER repair & Residual boundary\\
\midrule
\endfirsthead
\toprule
Difficulty & Invalid shortcut & FIBER repair & Residual boundary\\
\midrule
\endhead
Host truth versus object proof & Treat a Lean/Python theorem as \(\ipra\) derivability & Require literal local and final \texttt{Prf PRAAx}; compose by substitution and cut & Theorem-specific local derivations\\
Bounded quantifiers in induction & Induct on the quantified contract & Reflect to Boolean folds; induct on a quantifier-free state equation & Generic fold/reflection implementation\\
Hidden finite choice & Pick one witness for every node & Least bounded search plus master certificate, or seal local witnesses & Bound and decidable verifier\\
Nonunique witnesses & Concatenate arbitrary local witnesses into a run & Backward viability, finite strategy, or stable observation & Continuation and compatibility certificate\\
Cyclic precision demand & Iterate until host code happens to settle & Capped monotone closure with potential bound & Monotonicity and finite caps\\
Pointwise approximation & Call independent approximants a real name & Adjacent motion/nested intervals plus a tail modulus & Tail producer if not derivable\\
Metastability & Infer a global Cauchy rate & Keep a metastable interface; add summable tail for a name & Global rate is additional data\\
Exact sign/branch & Decide from arbitrary Cauchy names & Gap-guarded label or relation/cluster output & Positive separation\\
Repeated eigenvalue/identity rotation & Select eigenvector/axis globally & Projector, subspace, quaternion/projective observation & Spectral or distance guard for label\\
Division by nonzero & Use mere \(d\ne0\) numerically & Totalize and require \(|d|\ge\eta>0\) & Apartness margin\\
Abstract compactness & Select a minimizer or finite subcover & Finite net, coverage, modulus, local witness sealing & Effective net and modulus\\
Diophantine finiteness & Convert finiteness into a height cutoff & Bounded enumerator plus complete-height/list residual & Effective height or list theorem\\
Mordell--Weil/Selmer & Treat classes as rational representatives & Representative relation and finite viability & Representative totality\\
Ineffective constants & Store constants in host metadata & Put constants and thresholds in producer formulas & Effective estimate theorem\\
General probability space & Hide approximation and boundary mass & Finite rational laws, quantization, tails, anti-concentration & Quantitative approximation data\\
Completed topology & Reason directly with points of a completion & Finite complexes, nets, carriers, refinement observations & Approximation or separation producer\\
Code equality & Require canonical representatives & Observational setoids and congruence & Operation compatibility proofs\\
Finite gluing & Construct an infinitary quotient/sheaf & Local tuples, overlap checks, spanning-tree descent & Finite cover producer\\
Classical existence & Apply negative translation and claim a witness & Bounded strategies, reflected cases, or stable negative target & Explicit witness producer otherwise\\
Universal evaluator & Internalize all PR source programs & External finite specialization and generated proofs & Strong universal interface is impossible\\
\bottomrule
\end{longtable}
\setlength{\tabcolsep}{6pt}
\normalsize


\begin{thebibliography}{99}
\bibitem{BishopBridges}
E. Bishop and D. Bridges, \emph{Constructive Analysis}, Springer, 1985.

\bibitem{Kohlenbach}
U. Kohlenbach, \emph{Applied Proof Theory: Proof Interpretations and their Use in Mathematics}, Springer, 2008.

\bibitem{KohlenbachOliva}
U. Kohlenbach and P. Oliva, ``Proof mining: a systematic way of analysing proofs in mathematics,'' \emph{Proceedings of the Steklov Institute of Mathematics} 242 (2003), 136--164.

\bibitem{EscardoOliva}
M. Escard\'o and P. Oliva, ``Selection functions, bar recursion and backward induction,'' \emph{Mathematical Structures in Computer Science} 20 (2010), 127--168.

\bibitem{Salehi}
S. Salehi, ``Decidable formulas of intuitionistic primitive recursive arithmetic,'' \emph{Journal of Logic and Computation} 12 (2002), 101--114.

\bibitem{ProveItPRA}
V. Reshetnikov, \emph{ProveIt: Logic/PrimitiveRecursiveArithmetic}, repository snapshot at commit \texttt{58627ed89814ba122e1f0f012f1aba103cda267b}, 16 November 2026.  \url{https://github.com/VladimirReshetnikov/ProveIt/tree/main/Logic/PrimitiveRecursiveArithmetic}.

\bibitem{ConciseIPRA}
V. Reshetnikov, \emph{ConciseIPRA}, checked Lean presentation and bi-interpretation, in \cite{ProveItPRA}.

\bibitem{IPRADerivations}
V. Reshetnikov, \emph{IPRADerivations}, checked object-theory arithmetic library, in \cite{ProveItPRA}.

\bibitem{IPRABridge}
V. Reshetnikov, \emph{IPRABridge}, checked transport between the ordinary and Foundation presentations, in \cite{ProveItPRA}.
\end{thebibliography}

\end{document}
